%% Dokumentenklasse (Koma Script) -----------------------------------------
\documentclass[%
   %draft,     % Entwurfsstadium
   final,      % fertiges Dokument
	 % --- Paper Settings ---
   paper=a4,% [Todo: add alternatives]
   paper=portrait, % landscape
   pagesize=auto, % driver
   % --- Base Font Size ---
   fontsize=11pt,%
	 % --- Koma Script Version ---
   version=last, %
 ]{scrreprt} % Classes: scrartcl, scrreprt, scrbook
\usepackage{longtable}

% Encoding der Dateien (sonst funktionieren Umlaute nicht)
% Fuer Linux -> utf8
% Fuer Windows, alte Linux Distributionen -> latin1

% Empfohlen latin1, da einige Pakete mit utf8 Zeichen nicht
% funktionieren, z.B: listings, soul.
%\usepackage[latin1]{inputenc}
%\usepackage[ansinew]{inputenc}
\usepackage[utf8]{inputenc}
%\usepackage{ucs}
%\usepackage[utf8x]{inputenc}

%%% Preambel
%%% === Textbody ==============================================================
\KOMAoptions{%
   DIV=11,% (Size of Text Body, higher values = greater textbody)
   % DIV=calc % (also areaset/classic/current/default/last) 
   % -> after setting of spacing necessary!   
   BCOR=5mm% (Bindekorrektur)
}%
%%% === Headings ==============================================================
\KOMAoptions{%
   %%%% headings
	% headings=small,  % Small Font Size, thin spacing above and below
   % headings=normal, % Medium Font Size, medium spacing above and below
   headings=big, % Big Font Size, large spacing above and below
   %
   headings=noappendixprefix, % chapter in appendix as in body text
	headings=nochapterprefix,  % no prefix at chapters
   % headings=appendixprefix,   % inverse of 'noappendixprefix'
   % headings=chapterprefix,    % inverse of 'nochapterprefix'
   % headings=openany,   % Chapters start at any side
   % headings=openleft,  % Chapters start at left side
   headings=openright, % Chapters start at right side      
   %%% Add/Dont/Auto Dot behind section numbers 
   %%% (see DUDEN as reference)
   % numbers=autoenddot
   % numbers=enddot
   numbers=noenddot
   % secnumdepth=3 % depth of sections numbering (???)
}%
\setcounter{secnumdepth}{3}
%%% === Page Layout ===========================================================
\KOMAoptions{% (most options are for package typearea)
   twoside=true, % two side layout (alternating margins, standard in books)
   % twoside=false, % single side layout 
   % twoside=semi,  % two side layout (non alternating margins!)
   %
   twocolumn=false, % (true)
   %
   headinclude=false,%
   footinclude=false,%
   mpinclude=false,%      
   %
   headlines=2.1,%
	% headheight=2em,%
   headsepline=true,%
   footsepline=false,%
   cleardoublepage=empty %plain, headings
}%
%%% === Paragraph Separation ==================================================
\KOMAoptions{%
	% parskip=relative, % change indentation according to fontsize (recommanded)
   parskip=absolute, % do not change indentation according to fontsize
   parskip=false    % indentation of 1em
   % parskip=true   % parksip of 1 line - with free space in last line of 1em
   % parskip=full-  % parksip of 1 line - no adjustment
   % parskip=full+  % parksip of 1 line - with free space in last line of 1/4
   % parskip=full*  % parksip of 1 line - with free space in last line of 1/3
   % parskip=half   % parksip of 1/2 line - with free space in last line of 1em
   % parskip=half-  % parksip of 1/2 line - no adjustment
   % parskip=half+  % parksip of 1/2 line - with free space in last line of 1/3
   % parskip=half*  % parksip of 1/2 line - with free space in last line of 1em
}%
%%% === Table of Contents =====================================================
\setcounter{tocdepth}{3} % Depth of TOC Display
\KOMAoptions{%
   %%% Setting of 'Style' and 'Content' of TOC
   % toc=left, %
   toc=indented,%
   %
   toc=bib,
   % toc=nobib,
   % toc=bibnumbered,
   %
	% toc=index,%
   toc=noindex,
   %
   % toc=listof,
   toc=nolistof
   % toc=listofnumbered,
   %   
}%  
%%% === Lists of figures, tables etc. =========================================
\KOMAoptions{%
   %%% Setting of 'Style' and 'Content' of Lists 
   %%% (figures, tables etc)
	% --- General List Style ---
   listof=left, % tabular styles
   listof=indented, % hierarchical style
   % --- chapter highlighting ---
   % listof=chapterentry, % ??? Chapter starts are marked in figure/table
   % listof=chaptergapline, % New chapter starts are marked by a gap 
      		  			   	 % of a single line
	listof=chaptergapsmall, % New chapter starts are marked by a gap 
   	    					   % of a smallsingle line
   % listof=nochaptergap, % No Gap between chapters
   %
   % listof=leveldown, % lists are moved one level down ???
   % --- Appearance of Lists in TOC
   % listof=notoc, % Lists are not part of the TOC
   listof=totoc % add Lists to TOC without number
   % listof=totocnumbered, % add Lists to TOC with number
}%  
%%% === Bibliography ==========================================================
%% Setting of 'Style' and 'Content' of Bibliography
\KOMAoptions{%
	% bibliography=oldstyle,%
   bibliography=openstyle,%
   % bibliography=nottotoc, % Bibliography is not part of the TOC
   % bibliography=totocnumbered, % add Bibliography to TOC with number
   bibliography=totoc % add Bibliography to TOC without number
}%
%%% === Index =================================================================
%% Setting of 'Style' and 'Content' of Index in TOC
\KOMAoptions{%
   index=nottotoc % index is not part of the TOC
	% index=totoc, % add index to TOC without number
}%
%%% === Titlepage =============================================================
\KOMAoptions{%
   titlepage=true %
   %titlepage=false %
}%
%%% === Miscellaneous =========================================================
\KOMAoptions{% 	
   footnotes=multiple% nomultiple
   %open=any,%
   %open=left,%
   %open=right,%
   %chapterprefix=false,%
   %appendixprefix=false,%
   %chapteratlists=10pt,% entry
}%
% ------------------------------------------------------------------------
% LaTeX - Preambel  ******************************************************
% ------------------------------------------------------------------------
% von: Matthias Pospiech
% ========================================================================

% Strukturierung dieser Praeambel:
%    1.  Pakete die vor anderen geladen werden müssen
%        (calc, babel, xcolor, graphicx, amsmath, pst-pdf, ragged2e, ...)
%    2.  Schriften
%    3.  Mathematik (mathtools, fixmath, onlyamsmath, braket,
%        cancel, empheq, exscale, icomma, ...)
%    4.  Tabellen (booktabs, multirow, dcolumn, tabularx, ltxtable, supertabular)
%    5.  Text
%        5.1 Auszeichnungen (ulem, soul, url)
%        5.2 Fussnoten (footmisc)
%        5.3 Verweise (varioref)
%        5.4 Listen (enumitem, paralist, declist)
%    6.  Zitieren (csquotes, jurabib, natbib)
%    7.  PDF (microtype, hyperref, backref, hypcap, pdfpages
%    8.  Graphiken (float, flafter, placeins, subfig, wrapfig,
%        floatflt, picins, psfrag, sidecap, pict2e, curve2e)
%    9.  Sonstiges (makeidx, isodate, numprint, nomencl, acronym)
%    10. Verbatim (upquote, verbatim, fancyvrb, listings, examplep)
%    11. Wissenschaft (units)
%    12. Fancy Stuff
%    13. Layout
%       13.1.  Diverse Pakete und Einstellungen (multicol, ellipsis)
%       13.2.  Zeilenabstand (setspace)
%       13.3.  Seitenlayout (typearea, geometry)
%       13.4.  Farben
%       13.5.  Aussehen der URLS
%       13.6.  Kopf und Fusszeilen (scrpage2)
%       13.7.  Fussnoten
%       13.8.  Schriften (Sections )
%       13.9.  UeberSchriften (Chapter und Sections) (titlesec, indentfirst)
%       13.10. Captions (Schrift, Aussehen)
%              (caption, subfig, capt-of, mcaption, tocloft, multitoc, minitoc)
%    14.  Auszufuehrende Befehle
%    15. DoDo Notes Package


% ~~~~~~~~~~~~~~~~~~~~~~~~~~~~~~~~~~~~~~~~~~~~~~~~~~~~~~~~~~~~~~~~~~~~~~~~
% Einige Pakete muessen unbedingt vor allen anderen geladen werden
% ~~~~~~~~~~~~~~~~~~~~~~~~~~~~~~~~~~~~~~~~~~~~~~~~~~~~~~~~~~~~~~~~~~~~~~~~

%%% Packages for LaTeX - programming
%
% Define commands that don't eat spaces.
\usepackage{xspace}
% IfThenElse
\usepackage{ifthen}
%%% Doc: ftp://tug.ctan.org/pub/tex-archive/macros/latex/contrib/oberdiek/ifpdf.sty
% command for testing for pdf-creation
\usepackage{ifpdf} %\ifpdf \else \fi

%%% Internal Commands: ----------------------------------------------
\makeatletter
%
\providecommand{\IfPackageLoaded}[2]{\@ifpackageloaded{#1}{#2}{}}
\providecommand{\IfPackageNotLoaded}[2]{\@ifpackageloaded{#1}{}{#2}}
\providecommand{\IfElsePackageLoaded}[3]{\@ifpackageloaded{#1}{#2}{#3}}
%
\newboolean{chapteravailable}%
\setboolean{chapteravailable}{false}%

\ifcsname chapter\endcsname
  \setboolean{chapteravailable}{true}%
\else
  \setboolean{chapteravailable}{false}%
\fi


\providecommand{\IfChapterDefined}[1]{\ifthenelse{\boolean{chapteravailable}}{#1}{}}%
\providecommand{\IfElseChapterDefined}[2]{\ifthenelse{\boolean{chapteravailable}}{#1}{#2}}%

\providecommand{\IfDefined}[2]{%
\ifcsname #1\endcsname
   #2 %
\else
     % do nothing
\fi
}

\providecommand{\IfElseDefined}[3]{%
\ifcsname #1\endcsname
   #2 %
\else
   #3 %
\fi
}

\providecommand{\IfElseUnDefined}[3]{%
\ifcsname #1\endcsname
   #3 %
\else
   #2 %
\fi
}


%
% Check for 'draft' mode - commands.
\newcommand{\IfNotDraft}[1]{\ifx\@draft\@undefined #1 \fi}
\newcommand{\IfNotDraftElse}[2]{\ifx\@draft\@undefined #1 \else #2 \fi}
\newcommand{\IfDraft}[1]{\ifx\@draft\@undefined \else #1 \fi}
%

% Definde frontmatter, mainmatter and backmatter if not defined
\@ifundefined{frontmatter}{%
   \newcommand{\frontmatter}{%
      %In Roemischen Buchstaben nummerieren (i, ii, iii)
      \pagenumbering{roman}
   }
}{}
\@ifundefined{mainmatter}{%
   % scrpage2 benoetigt den folgenden switch
   % wenn \mainmatter definiert ist.
   \newif\if@mainmatter\@mainmattertrue
   \newcommand{\mainmatter}{%
      % -- Seitennummerierung auf Arabische Zahlen zuruecksetzen (1,2,3)
      \pagenumbering{arabic}%
      \setcounter{page}{1}%
   }
}{}
\@ifundefined{backmatter}{%
   \newcommand{\backmatter}{
      %In Roemischen Buchstaben nummerieren (i, ii, iii)
      \pagenumbering{roman}
   }
}{}

% Pakete speichern die spaeter geladen werden sollen
\newcommand{\LoadPackagesNow}{}
\newcommand{\LoadPackageLater}[1]{%
   \g@addto@macro{\LoadPackagesNow}{%
      \usepackage{#1}%
   }%
}



\makeatother
%%% ----------------------------------------------------------------
%
%%% Doc: www.cs.brown.edu/system/software/latex/doc/calc.pdf
% Calculation with LaTeX
\usepackage{calc}

%%% Doc: ftp://tug.ctan.org/pub/tex-archive/macros/latex/required/babel/babel.pdf
% Languagesetting
\usepackage[
%	german,
	ngerman,
%	english,
%	french,
]{babel}

%%% Doc: ftp://tug.ctan.org/pub/tex-archive/macros/latex/contrib/xcolor/xcolor.pdf
% Farben
% Incompatible: Do not load when using pstricks !
\usepackage[
	table % Load for using rowcolors command in tables
]{xcolor}


%%% Doc: ftp://tug.ctan.org/pub/tex-archive/macros/latex/required/graphics/grfguide.pdf
% Bilder
\usepackage[%
	%final,
	%draft % do not include images (faster)
]{graphicx}

%%% Doc: ftp://tug.ctan.org/pub/tex-archive/macros/latex/contrib/oberdiek/epstopdf.pdf
%% If an eps image is detected, epstopdf is automatically called to convert it to pdf format.
%% Requires: graphicx loaded
\usepackage{epstopdf}


%%% Doc: ftp://tug.ctan.org/pub/tex-archive/macros/latex/required/amslatex/math/amsldoc.pdf
% Amsmath - Mathematik Basispaket
%
% fuer pst-pdf displaymath Modus vor pst-pdf benoetigt.
\usepackage[
   centertags, % (default) center tags vertically
   %tbtags,    % 'Top-or-bottom tags': For a split equation, place equation numbers level
               % with the last (resp. first) line, if numbers are on the right (resp. left).
   sumlimits,  %(default) Place the subscripts and superscripts of summation
               % symbols above and below
   %nosumlimits, % Always place the subscripts and superscripts of summation-type
               % symbols to the side, even in displayed equations.
   intlimits,  % Like sumlimits, but for integral symbols.
   %nointlimits, % (default) Opposite of intlimits.
   namelimits, % (default) Like sumlimits, but for certain 'operator names' such as
               % det, inf, lim, max, min, that traditionally have subscripts placed underneath
               % when they occur in a displayed equation.
   %nonamelimits, % Opposite of namelimits.
   %leqno,     % Place equation numbers on the left.
   %reqno,     % Place equation numbers on the right.
   fleqn,     % Position equations at a fixed indent from the left margin
   			  % rather than centered in the text column.
]{amsmath} %
% eqnarray nicht zusammen mit amsmath benutzen, siehe l2tabu.pdf für
% Hintergruende.

%%% Doc: http://www.ctan.org/tex-archive/macros/latex/contrib/pst-pdf/pst-pdf-DE.pdf
% Used to automatically integrate eps graphics in an pdf document using pdflatex.
% Requires ps4pdf macro !!!
% Download macro from http://www.ctan.org/tex-archive/macros/latex/contrib/pst-pdf/scripts/
%
%\usepackage[%
%   %active,       % Aktiviert den Extraktionsmodus (DVI-Ausgabe). Die explizite Angabe ist
%                  % normalerweise unnötig (Standard im LATEX-Modus).
%   %inactive,     % Das Paket wird deaktiviert, Zuätzlich werden die Pakete pstricks und
%                  % graphicx geladen
%   nopstricks,    % Das Paket pstricks wird nicht geladen.
%   %draft,        % Im pdfLATEX-Modus werden aus der Containerdatei eingefügte Grafiken nur
%                  % als Rahmen dargestellt.
%   %final,        % Im pdfLATEX-Modus werden aus der Containerdatei eingefügte Grafiken
%                  % vollständig dargestellt (Standard).
%   %tightpage,    % Die Abmessung Grafiken in der Containerdatei entsprechen denen der
%                  % zugehörigen TEX-Boxen (Standard).
%   %notightpage,  % die Grafiken in der Containerdatei nehmen
%                  % mindestens die Größe des gesamten Blattes einnehmen.
%   %displaymath,  % Es werden zusätzlich die mathematischen Umgebungen displaymath,
%                  % eqnarray und $$ extrahiert und im pdf-Modus als Grafik eingefügt.
%]{pst-pdf}
%
% Notwendiger Bugfix für natbib Paket bei Benutzung von pst-pdf (Version <= v1.1o)
\IfPackageLoaded{pst-pdf}{
   \providecommand\makeindex{}
   \providecommand\makeglossary{}
}{}

%% Doc: ftp://tug.ctan.org/pub/tex-archive/graphics/pstricks/README
% load before graphicx
% \usepackage{pstricks}
% \usepackage{pst-plot, pst-node, pst-coil, pst-eps}

% This package implements a workaround for the LaTeX bug that marginpars
% sometimes appear on the wrong margin.
% \usepackage{mparhack}
% in some case this causes an error in the index together with package pdfpages
% the reason is unkown. Therefore I recommend to use the margins of marginnote

%% Doc: ftp://tug.ctan.org/pub/tex-archive/macros/latex/contrib/marginnote/marginnote.pdf
% Summary description: marginnote allows margin note, where \marginpar fails
\usepackage{marginnote}


%% Doc: (inside relsize.sty )
%% ftp://tug.ctan.org/pub/tex-archive/macros/latex/contrib/misc/relsize.sty
%  Set the font size relative to the current font size
\usepackage{relsize}

%% Doc: ftp://tug.ctan.org/pub/tex-archive/macros/latex/contrib/ms/ragged2e.pdf
% Besserer Flatternsatz (Linksbuendig, statt Blocksatz)
\usepackage{ragged2e}

% ~~~~~~~~~~~~~~~~~~~~~~~~~~~~~~~~~~~~~~~~~~~~~~~~~~~~~~~~~~~~~~~~~~~~~~~~
% Fonts Fonts Fonts
% ~~~~~~~~~~~~~~~~~~~~~~~~~~~~~~~~~~~~~~~~~~~~~~~~~~~~~~~~~~~~~~~~~~~~~~~~

\usepackage[T1]{fontenc} % T1 Schrift Encoding
\usepackage{textcomp}	 % Zusatzliche Symbole (Text Companion font extension)

%%% Schriften werden in Fonts.tex geladen
% ~~~~~~~~~~~~~~~~~~~~~~~~~~~~~~~~~~~~~~~~~~~~~~~~~~~~~~~~~~~~~~~~~~~~~~~~
% Fonts Fonts Fonts
% ~~~~~~~~~~~~~~~~~~~~~~~~~~~~~~~~~~~~~~~~~~~~~~~~~~~~~~~~~~~~~~~~~~~~~~~~

% Alle Schriften die hier angegeben sind sehen im PDF richtig aus.
% Die LaTeX Standardschrift ist die Latin Modern (lmodern Paket).
% If Latin Modern is not available for your distribution you must install the
% package cm-super instead. Otherwise your fonts will look horrible in the PDF

% DO NOT LOAD ae Package for the font !

%% ==== Zusammengesetzte Schriften  (Sans + Serif) =======================

%% - Latin Modern
\usepackage{lmodern}
%% -------------------
%
%% - Times, Helvetica, Courier (Word Standard...)
%\usepackage{mathptmx}
%\usepackage[scaled=.90]{helvet}
%\usepackage{courier}
%% -------------------
%%
%% - Palantino , Helvetica, Courier
%\usepackage{mathpazo}
%\usepackage[scaled=.95]{helvet}
%\usepackage{courier}
%% -------------------
%
%% - Bera Schriften
%\usepackage{bera}
%% -------------------
%
%% - Charter, Bera Sans
%\usepackage{charter}\linespread{1.05}
%\renewcommand{\sfdefault}{fvs}

%% ===== Serifen =========================================================

%\usepackage{mathpazo}                 %% --- Palantino
%\usepackage{charter}\linespread{1.05} %% --- Charter
%\usepackage{bookman}                  %% --- Bookman (laedt Avant Garde !!)
%\usepackage{newcent}                  %% --- New Century Schoolbook (laedt Avant Garde !!)

%\usepackage[%                         %% --- Fourier
%   upright,     % Math fonts are upright
%   expert,      % Only for EXPERT Fonts!
%   oldstyle,    % Only for EXPERT Fonts!
%   fulloldstyle % Only for EXPERT Fonts!
%]{fourier} %



%% ===== Sans Serif ======================================================

%\usepackage[scaled=.95]{helvet}        %% --- Helvetica
%\usepackage{cmbright}                  %% --- CM-Bright (eigntlich eine Familie)
%\usepackage{tpslifonts}                %% --- tpslifonts % Font for Slides
%\usepackage{avantgar}                  %% --- Avantgarde

%%%% =========== Italics ================

%\usepackage{chancery}                  %% --- Zapf Chancery

%%%% =========== Typewriter =============

%\usepackage{courier}                   %% --- Courier
%\renewcommand{\ttdefault}{cmtl}        %% --- CmBright Typewriter Font
%\usepackage[%                          %% --- Luxi Mono (Typewriter)
%   scaled=0.9
%]{luximono}



%%%% =========== Mathe ================

%% Recommanded to use with fonts: Aldus, Garamond, Melior, Sabon
%\usepackage[                           %% --- EulerVM (MATH)
%   small,       %for smaller Fonts
%  euler-digits % digits in euler fonts style
%]{eulervm}

% \usepackage[
% %   utopia,
% %   garamond,
%    charter
% ]{mathdesign}

%%%% (((( !!! kommerzielle Schriften !!! )))))))))))))))))))))))))))))))))))))))))))))))))))

%% ===== Serifen (kommerzielle Schriften ) ================================

%% --- Adobe Aldus
%\renewcommand{\rmdefault}{pasx}
%\renewcommand{\rmdefault}{pasj} %%oldstyle digits
% math recommended: \usepackage[small]{eulervm}

%% --- Adobe Garamond
%\usepackage[%
%   osf,        % oldstyle digits
%   scaled=1.05 %appropriate in many cases
%]{xagaramon}
% math recommended: \usepackage{eulervm}

%% --- Adobe Stempel Garamond
%\renewcommand{\rmdefault}{pegx}
%\renewcommand{\rmdefault}{pegj} %%oldstyle digits

%% --- Adobe Melior
%\renewcommand{\rmdefault}{pml}
% math recommended: %\usepackage{eulervm}

%% --- Adobe Minion
%\renewcommand{\rmdefault}{pmnx}
%\renewcommand{\rmdefault}{pmnj} %oldstyle digits
% math recommended: \usepackage[small]{eulervm} or \usepackage{mathpmnt} % commercial

%% --- Adobe Sabon
%\renewcommand{\rmdefault}{psbx}
%\renewcommand{\rmdefault}{psbj} %oldstyle digits
% math recommended: \usepackage{eulervm}

%% --- Adobe Times
% math recommended: \usepackage{mathptmx} % load first !
%\renewcommand{\rmdefault}{ptmx}
%\renewcommand{\rmdefault}{ptmj} %oldstyle digits

%% --- Linotype ITC Charter
%\renewcommand{\rmdefault}{lch}

%% --- Linotype Meridien
%\renewcommand{\rmdefault}{lmd}

%%% ===== Sans Serif (kommerzielle Schriften) ============================

%% --- Adobe Frutiger
%\usepackage[
%   scaled=0.90
%]{frutiger}

%% --- Adobe Futura (=Linotype FuturaLT) : Sans Serif
%\usepackage[
%   scaled=0.94  % appropriate in many cases
%]{futura}

%% --- Adobe Gill Sans : Sans Serif
%\usepackage{gillsans}

%% -- Adobe Myriad  : Sans Serif
%\renewcommand{\sfdefault}{pmy}
%\renewcommand{\sfdefault}{pmyc} %% condensed Font

%% --- Syntax : sans serif font
%\usepackage[
%   scaled
%]{asyntax}

%% --- Adobe Optima : Semi Sans Serif
%\usepackage[
%   medium %darker medium weight fonts
%]{optima}

%% --- Linotype ITC Officina Sans
%\renewcommand{\sfdefault}{lo9}





% ~~~~~~~~~~~~~~~~~~~~~~~~~~~~~~~~~~~~~~~~~~~~~~~~~~~~~~~~~~~~~~~~~~~~~~~~
% Math Packages
% ~~~~~~~~~~~~~~~~~~~~~~~~~~~~~~~~~~~~~~~~~~~~~~~~~~~~~~~~~~~~~~~~~~~~~~~~

% *** Mathematik **************************************
%
% amsmath schon vorher geladen da es vor pst-pdf geladen werden muss


%%% Doc: ftp://tug.ctan.org/pub/tex-archive/macros/latex/contrib/mh/doc/mathtools.pdf
% Erweitert amsmath und behebt einige Bugs
\usepackage[fixamsmath,disallowspaces]{mathtools}

%%% Doc: http://www.ctan.org/info?id=fixmath
% LaTeX's default style of typesetting mathematics does not comply
% with the International Standards ISO31-0:1992 to ISO31-13:1992
% which indicate that uppercase Greek letters always be typset
% upright, as opposed to italic (even though they usually
% represent variables) and allow for typsetting of variables in a
% boldface italic style (even though the required fonts are
% available). This package ensures that uppercase Greek be typeset
% in italic style, that upright $\Delta$ and $\Omega$ symbols are
% available through the commands \upDelta and \upOmega; and
% provides a new math alphabet \mathbold for boldface
% italic letters, including Greek.
\usepackage{fixmath}

%%% Doc: ftp://tug.ctan.org/pub/tex-archive/macros/latex/contrib/onlyamsmath/onlyamsmath.dvi
% Warnt bei Benutzung von Befehlen die mit amsmath inkompatibel sind.
\usepackage[
	all,
	warning
]{onlyamsmath}


%------------------------------------------------------

% -- Vektor fett darstellen -----------------
% \let\oldvec\vec
% \def\vec#1{{\boldsymbol{#1}}} %Fetter Vektor
% \newcommand{\ve}{\vec} %
% -------------------------------------------


%%% Doc: ftp://tug.ctan.org/pub/tex-archive/macros/latex/contrib/misc/braket.sty
\usepackage{braket}  % Quantenmechanik Bracket Schreibweise

%%% Doc: ftp://tug.ctan.org/pub/tex-archive/macros/latex/contrib/misc/cancel.sty
\usepackage{cancel}  % Durchstreichen

%%% Doc: ftp://tug.ctan.org/pub/tex-archive/macros/latex/contrib/mh/doc/empheq.pdf
\usepackage{empheq}  % Hervorheben

%%% Doc: ftp://tug.ctan.org/pub/tex-archive/info/math/voss/mathmode/Mathmode.pdf
%\usepackage{exscale} % Skaliert Mathe-Modus Ausgaben in allen Umgebungen richtig.

%%% Doc: ftp://tug.ctan.org/pub/tex-archive/macros/latex/contrib/was/icomma.dtx
% Erlaubt die Benutzung von Kommas im Mathematikmodus
\usepackage{icomma}


%%% Doc: http://www.ctex.org/documents/packages/special/units.pdf
% \usepackage[nice]{nicefrac}

%%% Tauschen von Epsilon und andere:
% \let\ORGvarrho=\varrho
% \let\varrho=\rho
% \let\rho=\ORGvarrho
%
\let\ORGvarepsilon=\varepsilon
\let\varepsilon=\epsilon
\let\epsilon=\ORGvarepsilon
%
% \let\ORGvartheta=\vartheta
% \let\vartheta=\theta
% \let\theta=\ORGvartheta
%
% \let\ORGvarphi=\varphi
% \let\varphi=\phi
% \let\phi=\ORGvarphi

% ~~~~~~~~~~~~~~~~~~~~~~~~~~~~~~~~~~~~~~~~~~~~~~~~~~~~~~~~~~~~~~~~~~~~~~~~
% Symbole
% ~~~~~~~~~~~~~~~~~~~~~~~~~~~~~~~~~~~~~~~~~~~~~~~~~~~~~~~~~~~~~~~~~~~~~~~~
%
%%% General Doc: http://www.ctan.org/tex-archive/info/symbols/comprehensive/symbols-a4.pdf
%
%% Symbole für Mathematiksatz
%\usepackage{mathrsfs} %% Schreibschriftbuchstaben für den Mathematiksatz (nur Großbuchstaben)
%\usepackage{dsfont}   %% Double Stroke Fonts
%\usepackage[mathcal]{euscript} %% adds euler mathcal font
%\usepackage{amssymb}
\usepackage[Symbolsmallscale]{upgreek} % upright symbols from euler package [Euler] or Adobe Symbols [Symbol]
\usepackage[upmu]{gensymb}             % Option upmu

%% Allgemeine Symbole
%\usepackage{wasysym}  %% Doc: http://www.ctan.org/tex-archive/macros/latex/contrib/wasysym/wasysym.pdf
%\usepackage{marvosym} %% Symbole aus der marvosym Schrift
%\usepackage{pifont}   %% ZapfDingbats


% ~~~~~~~~~~~~~~~~~~~~~~~~~~~~~~~~~~~~~~~~~~~~~~~~~~~~~~~~~~~~~~~~~~~~~~~~
% Tables (Tabular)
% ~~~~~~~~~~~~~~~~~~~~~~~~~~~~~~~~~~~~~~~~~~~~~~~~~~~~~~~~~~~~~~~~~~~~~~~~

% Basispaket fuer alle Tabellenfunktionen
% -> wird automatisch durch andere Pakete geladen
% \usepackage{array}
%
% bessere Abstaende innerhalb der Tabelle (Layout))
% -------------------------------------------------
%%% Doc: ftp://tug.ctan.org/pub/tex-archive/macros/latex/contrib/booktabs/booktabs.pdf
\usepackage{booktabs}
%
% Farbige Tabellen
% ----------------
% Das Paket colortbl wird inzwischen automatisch durch xcolor geladen
%
% Erweiterte Funktionen innerhalb von Tabellen
% --------------------------------------------
%%% Doc: ftp://tug.ctan.org/pub/tex-archive/macros/latex/contrib/multirow/multirow.sty
\usepackage{multirow} % Mehrfachspalten
%
%%% Doc: Documentation inside dtx Package
\usepackage{dcolumn}  % Ausrichtung an Komma oder Punkt

%%% Neue Tabellen-Umgebungen:
% ---------------------------
% Spalten automatischer Breite:
%%% Doc: Documentation inside dtx Package
% \usepackage{tabularx}
% -> nach hyperref Laden
% -> wird von ltxtable geladen
% \LoadPackageLater{tabularx}


% Tabellen ueber mehere Seiten
% ----------------------------
%%% Doc: ftp://tug.ctan.org/pub/tex-archive/macros/latex/contrib/carlisle/ltxtable.pdf
% \usepackage{ltxtable} % Longtable + tabularx
                        % (multi-page tables) + (auto-sized columns in a fixed width table)
% -> nach hyperref laden
\LoadPackageLater{ltxtable}


%%% Doc: ftp://tug.ctan.org/pub/tex-archive/macros/latex/contrib/supertabular/supertabular.pdf
%\usepackage{supertabular}


% ~~~~~~~~~~~~~~~~~~~~~~~~~~~~~~~~~~~~~~~~~~~~~~~~~~~~~~~~~~~~~~~~~~~~~~~~
% text related packages
% ~~~~~~~~~~~~~~~~~~~~~~~~~~~~~~~~~~~~~~~~~~~~~~~~~~~~~~~~~~~~~~~~~~~~~~~~

%%% Textverzierungen/Auszeichnungen ======================================
%
%%% Doc: ftp://tug.ctan.org/pub/tex-archive/macros/latex/contrib/misc/ulem.sty
\usepackage[normalem]{ulem}      % Zum Unterstreichen
%%% Doc: ftp://tug.ctan.org/pub/tex-archive/macros/latex/contrib/soul/soul.pdf
\usepackage{soul}		            % Unterstreichen, Sperren
%%% Doc: ftp://tug.ctan.org/pub/tex-archive/macros/latex/contrib/misc/url.sty
\usepackage{url} % Setzen von URLs. In Verbindung mit hyperref sind diese auch aktive Links.

%%% Fussnoten/Endnoten ===================================================
%
%%% Doc: ftp://tug.ctan.org/pub/tex-archive/macros/latex/contrib/footmisc/footmisc.pdf
%
\usepackage[
   bottom,      % Footnotes appear always on bottom. This is necessary
                % especially when floats are used
   stable,      % Make footnotes stable in section titles
   perpage,     % Reset on each page
   %para,       % Place footnotes side by side of in one paragraph.
   %side,       % Place footnotes in the margin
   ragged,      % Use RaggedRight
   %norule,     % suppress rule above footnotes
   multiple,    % rearrange multiple footnotes intelligent in the text.
   %symbol,     % use symbols instead of numbers
]{footmisc}

\renewcommand*{\multfootsep}{,\nobreakspace}

\deffootnote%
   [1em]% width of marker
   {1.5em}% indentation (general)
   {1em}% indentation (par)
   {\textsubscript{\thefootnotemark}}%


%% Einruecken der Fussnote einstellen
%\setlength\footnotemargin{10pt}

%--- footnote counter documentweit durchlaufend ------------------------------
%\usepackage{chngcntr}
%\counterwithout{footnote}{chapter}
%-----------------------------------------------------------------------------

%%% Doc: ftp://tug.ctan.org/pub/tex-archive/macros/latex/contrib/misc/endnotes.sty
%\usepackage{endnotes}
% From the Documentation:
% To turn all the footnotes in your documents into endnotes, say
%
%     \let\footnote=\endnote
%
%  in your preamble, and then add something like
%
%     \newpage
%     \begingroup
%     \parindent 0pt
%     \parskip 2ex
%     \def\enotesize{\normalsize}
%     \theendnotes
%     \endgroup
%
% as the last thing in your document.  (But \theendnotes all
% by itself will work.)

%%% Verweise =============================================================
%
%%% Doc: Documentation inside dtx File
\usepackage[ngerman]{varioref} % Intelligente Querverweise

%%% Listen ===============================================================
%
%
%%% Doc: ftp://tug.ctan.org/pub/tex-archive/macros/latex/contrib/paralist/paralist.pdf
% \usepackage{paralist}
%
%%% Doc: ftp://tug.ctan.org/pub/tex-archive/macros/latex/contrib/enumitem/enumitem.pdf
% Better than 'paralist' and 'enumerate' because it uses a keyvalue interface !
% Do not load together with enumerate.
\IfPackageNotLoaded{enumerate}{
	\usepackage{enumitem}
}
%
%%% Doc: ftp://tug.ctan.org/pub/tex-archive/macros/latex/contrib/ncctools/doc/desclist.pdf
% Improved description environment
%\usepackage{declist}


% ~~~~~~~~~~~~~~~~~~~~~~~~~~~~~~~~~~~~~~~~~~~~~~~~~~~~~~~~~~~~~~~~~~~~~~~~
% Pakete zum Zitieren
% ~~~~~~~~~~~~~~~~~~~~~~~~~~~~~~~~~~~~~~~~~~~~~~~~~~~~~~~~~~~~~~~~~~~~~~~~

% Quotes =================================================================
%% Doc: ftp://tug.ctan.org/pub/tex-archive/macros/latex/contrib/csquotes/csquotes.pdf
% Advanced features for clever quotations
\usepackage[%
   babel,            % the style of all quotation marks will be adapted
                     % to the document language as chosen by 'babel'
   german=quotes,		% Styles of quotes in each language
   english=british,
   french=guillemets
]{csquotes}

% All facilities which take a 'cite' argument will not insert
% it directly. They pass it to an auxiliary command called \mkcitation
% which  may be redefined to format the citation.
\renewcommand*{\mkcitation}[1]{{\,}#1}
\renewcommand*{\mkccitation}[1]{ #1}

\SetBlockThreshold{2} % Anzahl von Zeilen

\newenvironment{myquote}%
	{\begin{quote}\small}%
	{\end{quote}}%
\SetBlockEnvironment{myquote}
%\SetCiteCommand{} % Changes citation command


% Zitate =================================================================
%
% Reference by number
% \makeatletter
% \renewcommand\@biblabel[1]{#1.}
% \makeatother

% -------- Reference by author -------------------------------------------

% Doc: http://www.berger-on.net/jurabib/
% \usepackage{jurabib}
% \jurabibsetup{
%    authorformat={
%       %abbrv,%  First names will be abbreviated
%       %allreversed,% Names will be printed reversed ('first surname' in text and bibliography)
%       %citationreversed,% Names will be printed reversed ('first surname' only in text)
%       and,% Author separation with ',' and ', and' instead of the default slashes
%       %dynamic,% Font of author depends on existence of coauthor
%       %firstnotreversed, %All author names (except the first) will be printed reversed (in text only)
%       %indexed,% All author names are indexed separately (makeidx or index have to be loaded appropriately)
%       %italic,% Author will appear in italics
%       %reducedifibidem,%  Author names are reduced to the surname for subsequent
%       %citations,% (for the ibidem=name options)
%       smallcaps,% Author will appear in small caps
%       year,% Emulates author-year citations
%    },
%    bibformat={
%       %compress,% Reduces the vertical space between the items in the bibliography
%       %ibidem,%  Replaces repeated authors in the bibliography
%       %ibidemalt,% Special format for German law students
%       %nohang,% No hanging indent for bibliography
%       numbered,% Numbered items in bibliography
%       %raggedright,% Flushleft for bibliography
%       %tabular,% Tabular-like bibliography format
%    },
%    titleformat={
%       all,% Prints out all titles, doesn't care about multiple works
%       %colonsep,% Separation between author and title with colon
%       comma,% Separation between author and title with comma
%       italic,% Title will appear in italics
%    },
%    %biblikecite=true,% Formatting of bibliography follows formatting of citations (as far as possible)
%    %coauthorformat={
%    %  normal,% No special format for coauthors
%    %  %italic,% Coauthor will appear in italics
%    %},
%    %colastsep=divis,% Coauthor after author, separation with divis, Standard is a slash
%    %cofirstsep={
%    %% in,% Author after coauthor, separation with 'in'
%    %  comma,% Author after coauthor, separation with comma
%    %},
%    commabeforerest,  % Nach allen Angaben und vor den zusaetzlichen
%                      % (z.B. Seitenanzahl) wird ein Komma gesetzt
%    citefull={
%       %all,% All citations are full citations
%       first,% First citation is printed full
%       %chapter,% citefull=first, resetted each chapter (book and report classes)
%       %section,% citefull=first, resetted each section (article classes)
%    },
%    %chicago=true, %chicago-like format of citation and bibliography
%    %oxford=true, %Emulates oxford-like format of citation and bibliography
%    %crossref={
%    %  dynamic,% Long crossref's if they are used first time, shorter for all further citations
%    %  long,% Always long crossref's
%    %  short,% Always short crossref's (short as possible, longer if citations are ambiguous)
%    %},
%    %edby=true,% Switches from '(ed.)' to 'edited by' for incollections
%    %endnote=true,% The note field is printed at the end of the entry, after the closing period
%    %footnotes=marginal,% Another footnote format
%    %human=true, % Common humanities option, make authorformat=and the default
%    %howcited={
%    %  all,% The howcited remark is printed for all entries
%    %  compare,% The howcited remark is printed for works, where title and shorttitle differ
%    %  normal,% The howcited remark is printed for entries containing a non-empty howcited field
%    %  multiple,%  The howcited remark is printed if more than one work of the author is cited
%    %},
%    ibidem={
%       %name,% Ibidem with authors name
%       %name&title,% Ibidem with authors name and title
%       %nostrict,% Ibidem is allowed for every footnote
%       strict,% Ibidem is not allowed for first footnote on each page
%       %strictdoublepage,% Ibidem is not allowed for first footnote on left (even) pages
%    },
%    idem={
%       %nostrict,% Idem is allowed for every footnote
%       strict,% Idem is not allowed for first footnote on each page
%       %strictdoublepage,%  Idem is not allowed for first footnote on left (even) pages
%    },
%    %lookat=true,% Enables crossref to full (first) footnote citation
%    %natoptargorder=true,% Reversed optional arguments
%    %opcit={
%    %  true,% Enables op. cit. for already cited, but not subsequent works
%    %  chapter,% Resets opcit each chapter (book and report classes)
%    %  section,% Resets opcit each section (article classes)
%    %},
%    pages={
%       %always,%  Page(ranges)s given via the pages-field are always printed in the citation
%       format,%  Pages given via the optional argument and page(ranges)s given via the pages-field are formatted automatically
%       %test,% Page(ranges)s given via the pages-field are printed in the citation, if no pages are given by the optional argument
%    },
%    %see=true,% The second optional argument can be used to add sequences like 'see' before the citation
%    %superscriptedition={
%    %  all,% Superscripted edition number for all citations
%    %  bib,% Superscripted edition number for the bibliography
%    %  commented,% Superscripted edition number only for type @COMMENTED
%    %  switch,%  Superscripted edition number for works with field ssedition=1
%    %},
%    super, % alle \cite werden zu footcite
% }
%
% \bibliographystyle{jurabib} %
% %\bibliographystyle{jureco}
% %\bibliographystyle{jurunsrt}
% %\bibliographystyle{jox}
%
% \renewcommand{\biblnfont}{\bfseries\scshape\RaggedRight} % Autoren
% \renewcommand{\bibfnfont}{} % Autoren Vornamen
% \renewcommand{\bibelnfont}{\normalfont} % Herausgeber
% \renewcommand{\bibefnfont}{\normalfont} % Herausgeber Vornamen
% % Anpassung der Titel von Buechern
% \renewcommand{\bibtfont}{\normalfont\textit} % Titel
% % % Modifizierung des Zeitschriftentitels bei Artikeln.
% \renewcommand{\bibjtfont}{\normalfont}
% % % Titel eines Artikels, eines Beitrages in einem Sammelwerk oder aehnliches zu formatieren.
% \renewcommand{\bibapifont}{\normalfont\textit} % Periodical-Titel
% % % Aussehen des series Feldes bestimmen
% \renewcommand{\bibsnfont}{\textbf}
% %
% \renewcommand{\biburlprefix}{Webseite:{ }}
% \renewcommand{\biburlsuffix}{}
%
% % Ausgabe des Jahres
% \renewcommand{\jbcitationyearformat}[1]{(#1)}
%
% \renewcommand{\bibleftcolumnadjust}{\RaggedRight}
% \renewcommand{\bibrightcolumnadjust}{\RaggedRight}
% %
% -------- Reference by number (author) ----------------------------------

% %%% Doc: ftp://tug.ctan.org/pub/tex-archive/macros/latex/contrib/natbib/natbib.pdf
%\usepackage[%
%	%round,	%(default) for round parentheses;
%	square,	% for square brackets;
%	%curly,	% for curly braces;
%	%angle,	% for angle brackets;
%	%colon,	% (default) to separate multiple citations with colons;
%	comma,	% to use commas as separaters;
%	%authoryear,% (default) for author-year citations;
%	numbers,	% for numerical citations;
%	%super,	% for superscripted numerical citations, as in Nature;
%	sort,		% orders multiple citations into the sequence in which they appear in the list of references;
%	sort&compress,    % as sort but in addition multiple numerical citations
%                   % are compressed if possible (as 3-6, 15);
%	%longnamesfirst,  % makes the first citation of any reference the equivalent of
%                   % the starred variant (full author list) and subsequent citations
%                   %normal (abbreviated list);
%	%sectionbib,      % redefines \thebibliography to issue \section* instead of \chapter*;
%                   % valid only for classes with a \chapter command;
%                   % to be used with the chapterbib package;
%	%nonamebreak,     % keeps all the authors names in a citation on one line;
%                   %causes overfull hboxes but helps with some hyperref problems.
%]{natbib}

%\RequirePackage[fixlanguage]{babelbib}
%\bibliographystyle{babalpha-fl}%

%%% Bibliography styles with natbib support
%\bibliographystyle{plainnat} % Numeric Labels, alphabatical order
%\bibliographystyle{abbrvnat} % same as plain, but shorter names
%\bibliographystyle{unsrtnat} % same as plain, but appeariance in order of citation
%\bibliographystyle{alpha}    % labels are formed by author and year

%%% Bibliography styles according to DIN
%%% get from: http://www.ctan.org/tex-archive/biblio/bibtex/contrib/german/din1505/
%\bibliographystyle{alphadin}
%\bibliographystyle{abbrvdin}
%\bibliographystyle{plaindin}
%\bibliographystyle{unsrtdin}
%\bibliographystyle{bib/bst/alphadin-mod} % Modifiziert: Kleinere Abstaende vor ";" und kein "+" bei etal.

%%% Bibliography styles created with custombib
%%% Doc: ftp://tug.ctan.org/pub/tex-archive/macros/latex/contrib/custom-bib/makebst.pdf
\bibliographystyle{bib/bst/AlphaDINFirstName}

% other BibTeX styles: http://www.cs.stir.ac.uk/~kjt/software/latex/showbst.html


% ~~~~~~~~~~~~~~~~~~~~~~~~~~~~~~~~~~~~~~~~~~~~~~~~~~~~~~~~~~~~~~~~~~~~~~~~
% figures and placement
% ~~~~~~~~~~~~~~~~~~~~~~~~~~~~~~~~~~~~~~~~~~~~~~~~~~~~~~~~~~~~~~~~~~~~~~~~

%% Bilder und Graphiken ==================================================

%%% Doc: only dtx Package
\usepackage{float}             % Stellt die Option [H] fuer Floats zur Verfgung

%%% Doc: No Documentation
\usepackage{flafter}          % Floats immer erst nach der Referenz setzen

% Defines a \FloatBarrier command, beyond which floats may not
% pass; useful, for example, to ensure all floats for a section
% appear before the next \section command.
\usepackage[
	section		% "\section" command will be redefined with "\FloatBarrier"
]{placeins}
%
%%% Doc: ftp://tug.ctan.org/pub/tex-archive/macros/latex/contrib/subfig/subfig.pdf
% Incompatible: loads package capt-of. Loading of 'capt-of' afterwards will fail therefor
\usepackage{subfig} % Layout wird weiter unten festgelegt !

%%% Bilder von Text Umfliessen lassen : (empfehle wrapfig)
%
%%% Doc: ftp://tug.ctan.org/pub/tex-archive/macros/latex/contrib/wrapfig/wrapfig.sty
\usepackage{wrapfig}	        % defines wrapfigure and wrapfloat
%\setlength{\wrapoverhang}{\marginparwidth} % aeerlapp des Bildes ...
%\addtolength{\wrapoverhang}{\marginparsep} % ... in den margin
\setlength{\intextsep}{0.5\baselineskip} % Platz ober- und unterhalb des Bildes
% \intextsep ignoiert bei draft ???
%\setlength{\columnsep}{1em} % Abstand zum Text

%%% Doc: Documentation inside dtx Package
%\usepackage{floatflt}   	  % LaTeX2e Paket von 1996
                             % [rflt] - Standard float auf der rechten Seite

%%% Doc: ftp://tug.ctan.org/pub/tex-archive/macros/latex209/contrib/picins/picins.doc
%\usepackage{picins}          % LaTeX 2.09 Paket von 1992. aber Layout kombatibel


% Make float placement easier
\renewcommand{\floatpagefraction}{.75} % vorher: .5
\renewcommand{\textfraction}{.1}       % vorher: .2
\renewcommand{\topfraction}{.8}        % vorher: .7
\renewcommand{\bottomfraction}{.5}     % vorher: .3
\setcounter{topnumber}{3}              % vorher: 2
\setcounter{bottomnumber}{2}           % vorher: 1
\setcounter{totalnumber}{5}            % vorher: 3


%%% Doc: ftp://tug.ctan.org/pub/tex-archive/macros/latex/contrib/psfrag/pfgguide.pdf
% \usepackage{psfrag}	% Ersetzen von Zeichen in eps Bildern


%%% Doc: http://www.ctan.org/tex-archive/macros/latex/contrib/sidecap/sidecap.pdf
\usepackage[%
%	outercaption,%	(default) caption is placed always on the outside side
%	innercaption,% caption placed on the inner side
%	leftcaption,%  caption placed on the left side
	rightcaption,% caption placed on the right side
%	wide,%			caption of float my extend into the margin if necessary
%	margincaption,% caption set into margin
	ragged,% caption is set ragged
]{sidecap}

\renewcommand\sidecaptionsep{2em}
%\renewcommand\sidecaptionrelwidth{20}
\sidecaptionvpos{table}{c}
\sidecaptionvpos{figure}{c}


%% Diagramme mit LaTeX ===================================================
%

%%% Doc: ftp://tug.ctan.org/pub/tex-archive/macros/latex/contrib/pict2e/pict2e.pdf
% Neuimplementation der Picture Umgebung.
%
% The new package extends the existing LaTeX picture environment, using
% the familiar technique (cf. the graphics and color packages) of driver
% files.  The package documentation (pict2e.dtx) has a fair number of
% examples of use, showing where things are improved by comparison with
% the LaTeX picture environment.
% \usepackage{pict2e}

%%% Doc: ftp://tug.ctan.org/pub/tex-archive/macros/latex/contrib/curve2e/curve2e.pdf
% Extensions for package pict2e.
%\usepackage{curve2e}
%


% ~~~~~~~~~~~~~~~~~~~~~~~~~~~~~~~~~~~~~~~~~~~~~~~~~~~~~~~~~~~~~~~~~~~~~~~~
% misc packages
% ~~~~~~~~~~~~~~~~~~~~~~~~~~~~~~~~~~~~~~~~~~~~~~~~~~~~~~~~~~~~~~~~~~~~~~~~

\usepackage{makeidx}		% Index
\IfDraft{
  \usepackage{showidx}    % Indizierte Begriffe am Rand (Korrekturlesen)
}

%%% Doc: ftp://tug.ctan.org/pub/tex-archive/macros/latex/contrib/isodate/README
%%% Incompatible: draftcopy
% Tune the output format of dates.
%\usepackage{isodate}

%%% Doc: ftp://tug.ctan.org/pub/tex-archive/macros/latex/contrib/numprint/numprint.pdf
% Modify printing of numbers
%\usepackage{numprint}

%%% Doc: ftp://tug.ctan.org/pub/tex-archive/macros/latex/contrib/nomencl/nomencl.pdf
\usepackage[%
	german,
	%english
]{nomencl}[2005/09/22]

\usepackage[
	footnote,	% Full names appear in the footnote
	smaller,		% Print acronym in smaller fontsize
	printonlyused %
]{acronym}

% ~~~~~~~~~~~~~~~~~~~~~~~~~~~~~~~~~~~~~~~~~~~~~~~~~~~~~~~~~~~~~~~~~~~~~~~~
% verbatim packages
% ~~~~~~~~~~~~~~~~~~~~~~~~~~~~~~~~~~~~~~~~~~~~~~~~~~~~~~~~~~~~~~~~~~~~~~~~

%%% Doc: ftp://tug.ctan.org/pub/tex-archive/macros/latex/contrib/upquote/upquote.sty
\usepackage{upquote} % Setzt "richtige" Quotes in verbatim-Umgebung

%%% Doc: No Documentation
% \usepackage{verbatim} %Reimplemntation of the original verbatim

%%% Doc: http://www.cs.brown.edu/system/software/latex/doc/fancyvrb.pdf
% \usepackage{fancyvrb} % Superior Verbatim Class

%% Listings Paket ------------------------------------------------------

\definecolor{keywordblue}{RGB}{0,0,0}
\definecolor{stringred}{RGB}{0,0,0}
%\definecolor{keywordblue}{RGB}{127,0,127}
%\definecolor{stringred}{RGB}{196,26,22}
 \usepackage{listings}
	\lstset{basicstyle={\footnotesize\ttfamily},
	breaklines=true,
	extendedchars=true,
	frame=b,
	framexbottommargin=4pt,
	framexleftmargin=17pt,
	framexrightmargin=5pt,
	keywordstyle={\color{keywordblue}},
	numbers=left,
	numbersep=5pt,
	numberstyle={\tiny},
	showspaces=false,
	showstringspaces=false,
	showtabs=false,
	stringstyle={\color{stringred}\ttfamily},
	tabsize=2,
	xleftmargin=17pt}
\usepackage{courier}
 \lstloadlanguages{% Check Dokumentation for further languages ...
%         [Visual]Basic
%         %Pascal
%         %C
%         %C++
%         %XML
          HTML,
          PHP,
          XML,
          SQL
 }
 
\usepackage{caption}
 \DeclareCaptionFont{white}{\color{white}}
  \DeclareCaptionFormat{listing}{\colorbox[cmyk]{0.43, 0.35, 0.35,0.01}{\parbox{\textwidth}{\hspace{15pt}#1#2#3}}}
 \captionsetup[lstlisting]{format=listing,labelfont=white,textfont=white, singlelinecheck=false, margin=0pt, font={bf,footnotesize}}

%%% Doc: ftp://tug.ctan.org/pub/tex-archive/macros/latex/contrib/examplep/eurotex_2005_examplep.pdf
% LaTeX Code und Ergebnis nebeneinander darstellen
%\usepackage{examplep}

% ~~~~~~~~~~~~~~~~~~~~~~~~~~~~~~~~~~~~~~~~~~~~~~~~~~~~~~~~~~~~~~~~~~~~~~~~
% science packages
% ~~~~~~~~~~~~~~~~~~~~~~~~~~~~~~~~~~~~~~~~~~~~~~~~~~~~~~~~~~~~~~~~~~~~~~~~

\usepackage{units}


% ~~~~~~~~~~~~~~~~~~~~~~~~~~~~~~~~~~~~~~~~~~~~~~~~~~~~~~~~~~~~~~~~~~~~~~~~
% fancy packages
% ~~~~~~~~~~~~~~~~~~~~~~~~~~~~~~~~~~~~~~~~~~~~~~~~~~~~~~~~~~~~~~~~~~~~~~~~

%%% Doc: No documentation - documented in 'The LaTeX Companion'
% \usepackage{fancybox}   % for shadowbox, ovalbox

%%% Doc: ftp://tug.ctan.org/pub/tex-archive/macros/latex/contrib/misc/framed.sty
% \usepackage{framed}
% \renewcommand\FrameCommand{\fcolorbox{black}{shadecolor}}

\makeatletter
\IfPackageLoaded{framed}{%
   \IfPackageLoaded{marginnote}{%
      \begingroup
         \g@addto@macro\framed{%
            \let\marginnoteleftadjust\FrameSep
               \let\marginnoterightadjust\FrameSep
         }
      \makeatother
  }
}
\makeatother



%%% Doc: No documentation - documented in 'The LaTeX Companion'
% \usepackage{boxedminipage}

%%% Doc: ftp://tug.ctan.org/pub/tex-archive/macros/latex/contrib/lettrine/doc/lettrine.pdf
% Dropping capitals
% \usepackage{lettrine}


% ~~~~~~~~~~~~~~~~~~~~~~~~~~~~~~~~~~~~~~~~~~~~~~~~~~~~~~~~~~~~~~~~~~~~~~~~
% layout packages
% ~~~~~~~~~~~~~~~~~~~~~~~~~~~~~~~~~~~~~~~~~~~~~~~~~~~~~~~~~~~~~~~~~~~~~~~~

%%% Diverse Pakete und Einstellungen =====================================

%%% Doc: Documentation inside dtx file
% Mehere Text-Spalten
\usepackage{multicol}

%\nonfrenchspacing     % liefert extra Platz hinter Satzenden.
                       % Fuer deutschen Text standardmaessig ausgeschaltet!


\usepackage{ellipsis}  % >>Intelligente<< \dots

%% Zeilenabstand =========================================================
%
%%% Doc: ftp://tug.ctan.org/pub/tex-archive/macros/latex/contrib/setspace/setspace.sty
\usepackage{setspace}
\onehalfspacing		% 1,5-facher Abstand
%\doublespacing		% 2-facher Abstand
% hereafter load 'typearea' again


%% Hurenkinder und Schusterjungen unterdrücken
\clubpenalty = 10000  
\widowpenalty = 10000  
\displaywidowpenalty = 10000 

%% Seitenlayout ==========================================================
%
% Layout laden um im Dokument den Befehl \layout nutzen zu koennen
%%% Doc: no documentation
%\usepackage[verbose]{layout}
%

% Layout mit 'geometry'
%%% Doc: ftp://tug.ctan.org/pub/tex-archive/macros/latex/contrib/geometry/manual.pdf
\usepackage{geometry}

\IfPackageLoaded{geometry}{%
\geometry{%
%%% Paper Groesse
   a4paper, % Andere a0paper, a1paper, a2paper, a3paper, , a5paper, a6paper,
            % b0paper, b1paper, b2paper, b3paper, b4paper, b5paper, b6paper
            % letterpaper, executivepaper, legalpaper
   %screen,  % a special paper size with (W,H) = (225mm,180mm)
   %paperwidth=,
   %paperheight=,
   %papersize=, %{ width , height }
   %landscape,  % Querformat
   portrait,    % Hochformat
%%% Koerper Groesse
   %hscale=,      % ratio of width of total body to \paperwidth
                  % hscale=0.8 is equivalent to width=0.8\paperwidth. (0.7 by default)
   %vscale=,      % ratio of height of total body to \paperheight
                  % vscale=0.9 is equivalent to height=0.9\paperheight.
   %scale=,       % ratio of total body to the paper. scale={ h-scale , v-scale }
   %totalwidth=,    % width of total body % (Generally, width >= textwidth)
   %totalheight=,   % height of total body, excluding header and footer by default
   %total=,        % total={ width , height }
   %textwidth=,    % modifies \textwidth, the width of body
   %textheight=,   % modifies \textheight, the height of body
   %body=,        % { width , height } sets both \textwidth and \textheight of the body of page.
   %lines=45,       % enables users to specify \textheight by the number of lines.
   %includehead,  % includes the head of the page, \headheight and \headsep, into total body.
   %includefoot,  % includes the foot of the page, \footskip, into body.
   %includeheadfoot, % sets both includehead and includefoot to true
   %includemp,    % includes the margin notes, \marginparwidth and \marginparsep, into body
   %includeall,   % sets both includeheadfoot and includemp to true.
   %ignorehead,   % disregards the head of the page, headheight and headsep in determining vertical layout
   %ignorefoot,   % disregards the foot of page, footskip, in determining vertical layout
   %ignoreheadfoot, % sets both ignorehead and ignorefoot to true.
   %ignoremp,     % disregards the marginal notes in determining the horizontal margins
   ignoreall,     % sets both ignoreheadfoot and ignoremp to true
   heightrounded, % This option rounds \textheight to n-times (n: an integer) of \baselineskip
   %hdivide=,     % { left margin , width , right margin }
                  % Note that you should not specify all of the three parameters
   %vdivide=,     % { top margin , height , bottom margin }
   %divide=,      % ={A,B,C} %  is interpreted as hdivide={A,B,C} and vdivide={A,B,C}.
%%% Margin
   %left=,        % left margin (for oneside) or inner margin (for twoside) of total body
                  % alias: lmargin, inner
   %right=,       % right or outer margin of total body
                  % alias: rmargin outer
   %top=3cm,         % top margin of the page.
                  % Alias : tmargin
   %bottom=3cm,      % bottom margin of the page
                  % Alias : bmargin
   hmargin={3.5cm, 2.5cm },     % left and right margin. hmargin={ left margin , right margin }
   vmargin={3cm, 3cm },     % top and bottom margin. vmargin={ top margin , bottom margin }
   %margin=,      % margin={A,B} is equivalent to hmargin={A,B} and vmargin={A,B}
   %hmarginratio, % horizontal margin ratio of left (inner) to right (outer).
   %vmarginratio, % vertical margin ratio of top to bottom.
   %marginratio,  % marginratio={ horizontal ratio , vertical ratio }
   %hcentering,   % sets auto-centering horizontally and is equivalent to hmarginratio=1:1
   %vcentering,   % sets auto-centering vertically and is equivalent to vmarginratio=1:1
   %centering,    % sets auto-centering and is equivalent to marginratio=1:1
   %twoside,       % switches on twoside mode with left and right margins swapped on verso pages.
   %asymmetric,   % implements a twosided layout in which margins are not swapped on alternate pages
                  % and in which the marginal notes stay always on the same side.
   bindingoffset=5mm,  % removes a specified space for binding
%%% Dimensionen
   %headheight=,  % Alias:  head
   %headsep=,     % separation between header and text
   %footskip=,    % distance separation between baseline of last line of text and baseline of footer
                  % Alias: foot
   %nohead,       % eliminates spaces for the head of the page
                  % equivalent to both \headheight=0pt and \headsep=0pt.
   %nofoot,       % eliminates spaces for the foot of the page
                  % equivalent to \footskip=0pt.
   %noheadfoot,   % equivalent to nohead and nofoot.
   %footnotesep=, % changes the dimension \skip\footins,.
                  % separation between the bottom of text body and the top of footnote text
   marginparwidth=0pt, % width of the marginal notes
                  % Alias: marginpar
   %marginparsep=,% separation between body and marginal notes.
   %nomarginpar,  % shrinks spaces for marginal notes to 0pt
   %columnsep=,   % the separation between two columns in twocolumn mode.
   %hoffset=,
   %voffset=,
   %offset=,      % horizontal and vertical offset.
                  % offset={ hoffset , voffset }
   %twocolumn,    % twocolumn=false denotes onecolumn
   twoside,
   %textwidth=400pt,   % sets \textwidth directly
   %textheight=,  % sets \textheight directly
   %reversemp,    % makes the marginal notes appear in the left (inner) margin
                  % Alias: reversemarginpar
}
} % Endif

% - Anzeigen des Layouts -
\IfPackageLoaded{geometry}{%
   %\geometry{showframe}
}

% Layout mit 'typearea'
%%% Doc: ftp://tug.ctan.org/pub/tex-archive/macros/latex/contrib/koma-script/scrguide.pdf

\IfPackageLoaded{typearea}{% Wenn typearea geladen ist
   \IfPackageNotLoaded{geometry}{% aber nicht geometry
      \typearea[current]{last}
   }
}

% BCOR
%    current  % Satzspiegelberechnung mit dem aktuell gültigen BCOR-Wert erneut
%             % durchführen.
% DIV
%    calc     % Satzspiegelberechnung einschließlich Ermittlung eines guten
%             % DIV-Wertes erneut durchführen.
%    classic  % Satzspiegelberechnung nach dem
%             % mittelalterlichen Buchseitenkanon
%             % (Kreisberechnung) erneut durchführen.
%    current  % Satzspiegelberechnung mit dem aktuell gültigen DIV-Wert erneut
%             % durchführen.
%    default  % Satzspiegelberechnung mit dem Standardwert für das aktuelle
%             % Seitenformat und die aktuelle Schriftgröße erneut durchführen.
%             % Falls kein Standardwert existiert calc anwenden.
%    last     % Satzspiegelberechnung mit demselben DIV -Argument, das beim
%             % letzten Aufruf angegeben wurde, erneut durchführen


%\usepackage[colorgrid,texcoord,gridunit=mm]{showframe}

\raggedbottom     % Variable Seitenhoehen zulassen

% Farben ================================================================

\IfDefined{definecolor}{%

% Farbe der Ueberschriften
%\definecolor{sectioncolor}{RGB}{0, 51, 153} % Blau
%\definecolor{sectioncolor}{RGB}{0, 25, 152}    % Blau (dunkler))
\definecolor{sectioncolor}{RGB}{0, 0, 0}    % Schwarz
%
% Farbe des Textes
\definecolor{textcolor}{RGB}{0, 0, 0}        % Schwarz
%
% Farbe fuer grau hinterlegte Boxen (fuer Paket framed.sty)
\definecolor{shadecolor}{gray}{0.90}

% Farben fuer die Links im PDF
\definecolor{pdfurlcolor}{rgb}{0,0,0.6}
\definecolor{pdffilecolor}{rgb}{0.7,0,0}
\definecolor{pdflinkcolor}{rgb}{0,0,0.6}
\definecolor{pdfcitecolor}{rgb}{0,0,0.6}

%% PDF-Linkfarben auf schwarz für den Druck:
% \definecolor{pdfurlcolor}{rgb}{0,0,0}
% \definecolor{pdffilecolor}{rgb}{0,0,0}
% \definecolor{pdflinkcolor}{rgb}{0,0,0}
% \definecolor{pdfcitecolor}{rgb}{0,0,0}


% Farben fuer Listings
\colorlet{stringcolor}{green!40!black!100}
\colorlet{commencolor}{blue!0!black!100}

} % Endif

%% Aussehen der URLS======================================================

%fuer URL (nur wenn url geladen ist)
\IfDefined{urlstyle}{
	\urlstyle{tt} %sf
}

%% Kopf und Fusszeilen====================================================
%%% Doc: ftp://tug.ctan.org/pub/tex-archive/macros/latex/contrib/koma-script/scrguide.pdf

\usepackage[%
   % headtopline,
   % plainheadtopline,
   % headsepline,
   % plainheadsepline,
   % footsepline,
   % plainfootsepline,
   % footbotline,
   % plainfootbotline,
   % ilines,
   % clines,
   % olines,
   automark,
   % autooneside,% ignore optional argument in automark at oneside
   komastyle,
   % standardstyle,
   % markuppercase,
   % markusedcase,
   nouppercase,
]{scrpage2}

%\usepackage[%
%   automark,         % automatische Aktualisierung der Kolumnentitel
%   nouppercase,      % Grossbuchstaben verhindern
%   %markuppercase    % Grossbuchstaben erzwingen
%   %markusedcase     % vordefinierten Stil beibehalten
%   %komastyle,       % Stil von Koma Script
%   %standardstyle,   % Stil der Standardklassen
%]{scrpage2}

\IfElseChapterDefined{%
   \pagestyle{scrheadings} % Seite mit Headern
}{
   \pagestyle{scrplain} % Seiten ohne Header
}
%\pagestyle{empty} % Seiten ohne Header
%
% loescht voreingestellte Stile
\clearscrheadings
\clearscrplain
%
% Was steht wo...
\IfElseChapterDefined{
   % Oben aussen: Kapitel und Section
   % Unten aussen: Seitenzahl
     \ohead{\headmark} % Oben außen: Setzt Kapitel und Section automatisch
     \ofoot[\pagemark]{\pagemark}
   % oder...
   % Oben aussen: Seitenzahlen
   % Oben innen: Kapitel und Section
   % \ohead{\pagemark}
   % \ihead{\headmark}
   % \ofoot[\pagemark]{} % Außen unten: Seitenzahlen bei plain
}{
   \cfoot[\pagemark]{\pagemark} % Mitte unten: Seitenzahlen bei plain
}
% Vollstaendige Liste der moeglichen Positionierungen
% \lehead[scrplain-links-gerade]{scrheadings-links-gerade}
% \cehead[scrplain-mittig-gerade]{scrheadings-mittig-gerade}
% \rehead[scrplain-rechts-gerade]{scrheadings-rechts-gerade}
% \lefoot[scrplain-links-gerade]{scrheadings-links-gerade}
% \cefoot[scrplain-mittig-gerade]{scrheadings-mittig-gerade}
% \refoot[scrplain-rechts-gerade]{scrheadings-rechts-gerade}
% \lohead[scrplain-links-ungerade]{scrheadings-links-ungerade}
% \cohead[scrplain-mittig-ungerade]{scrheadings-mittig-ungerade}
% \rohead[scrplain-rechts-ungerade]{scrheadings-rechts-ungerade}
% \lofoot[scrplain-links-ungerade]{scrheadings-links-ungerade}
% \cofoot[scrplain-mittig-ungerade]{scrheadings-mittig-ungerade}
% \rofoot[scrplain-rechts-ungerade]{scrheadings-rechts-ungerade}
% \ihead[scrplain-innen]{scrheadings-innen}
% \chead[scrplain-zentriert]{scrheadings-zentriert}
% \ohead[scrplain-außen]{scrheadings-außen}
% \ifoot[scrplain-innen]{scrheadings-innen}
% \cfoot[scrplain-zentriert]{scrheadings-zentriert}
% \ofoot[scrplain-außen]{scrheadings-außen}


%\usepackage{lastpage} % Stellt 'LastPage' zur Verfuegung
%\cfoot[Seite \pagemark~von \pageref{LastPage}]{} % Seitenzahl von Anzahl Seiten

% Angezeigte Abschnitte im Header
\IfElseChapterDefined{
   \automark[section]{chapter} %[rechts]{links}
}{
   \automark[subsection]{section} %[rechts]{links}
}
%
% Linien (moegliche Kombination mit Breiten)
\IfChapterDefined{
   %\setheadtopline{}     % modifiziert die Parameter fuer die Linie ueber dem
   							  %	 Seitenkopf
   \setheadsepline{.4pt}[\color{black}]
                         % modifiziert die Parameter fuer die Linie zwischen
                         % Kopf und Textkörper
   %\setfootsepline{}    % modifiziert die Parameter fuer die Linie zwischen
   							 % Text und Fuß
   %\setfootbotline{}    % modifiziert die Parameter fuer die Linie unter dem
   							 % Seitenfuss
}



% Groesse des Headers

% Breite von Kopf und Fusszeile einstellen
% \setheadwidth[Verschiebung]{Breite}
% \setfootwidth[Verschiebung]{Breite}
% mögliche Werte
% paper - die Breite des Papiers
% page - die Breite der Seite
% text - die Breite des Textbereichs
% textwithmarginpar - die Breite des Textbereichs inklusive dem Seitenrand
% head - die aktuelle Breite des Seitenkopfes
% foot - die aktuelle Breite des Seitenfusses
\setheadwidth[0pt]{text}
\setfootwidth[0pt]{text}

%\setheadwidth[]%
%{%
%   paper % width of paper
%   page  % width of page (paper - BCOR)
%   text  % \textwidth
%   textwithmarginpar % width of text plus margin
%   head  % current width of head
%   foot  % current width of foot
%}%


%% Fussnoten =============================================================
% Keine hochgestellten Ziffern in der Fussnote (KOMA-Script-spezifisch):
\deffootnote{1.5em}{1em}{\makebox[1.5em][l]{\thefootnotemark}}
\addtolength{\skip\footins}{\baselineskip} % Abstand Text <-> Fussnote

\setlength{\dimen\footins}{10\baselineskip} % Beschraenkt den Platz von Fussnoten auf 10 Zeilen

\interfootnotelinepenalty=10000 % Verhindert das Fortsetzen von
                                % Fussnoten auf der gegenüberligenden Seite


%% Schriften (Sections )==================================================

\IfElsePackageLoaded{fourier}{
   \newcommand\SectionFontStyle{\rmfamily}
}{
   \newcommand\SectionFontStyle{\sffamily}
}

% -- Koma Schriften --
\IfChapterDefined{%
   \setkomafont{chapter}{\huge\SectionFontStyle}    % Chapter
}

\setkomafont{sectioning}{\SectionFontStyle}
%\setkomafont{part}{\usekomafont{sectioning}}
%\setkomafont{section}{\usekomafont{sectioning}}
%\setkomafont{subsection}{\usekomafont{sectioning}}
%\setkomafont{subsubsection}{\usekomafont{sectioning}}
%\setkomafont{paragraph}{\usekomafont{sectioning}}
%\setkomafont{subparagraph}{\usekomafont{sectioning}}

\setkomafont{descriptionlabel}{\itshape}

%\setkomafont{caption}{\normalfont}
%\setkomafont{captionlabel}{\normalfont}

%\setkomafont{dictum}{}
%\setkomafont{dictumauthor}{}
%\setkomafont{dictumtext}{}
%\setkomafont{disposition}{}
%\setkomafont{footnote}{}
%\setkomafont{footnotelabel}{}
%\setkomafont{footnotereference}{}
%\setkomafont{minisec}{}

%\setkomafont{partnumber}{\bfseries\SectionFontStyle}
%\setkomafont{partentrynumber}{}
%\setkomafont{chapterentrypagenumber}{}
%\setkomafont{sectionentrypagenumber}{}

\setkomafont{pageheadfoot}{\normalfont\normalcolor\small\sffamily}
\setkomafont{pagenumber}{\bfseries\usekomafont{sectioning}}

%\setkomafont{partentry}{\usekomafont{sectioning}\large}
%\setkomafont{chapterentry}{\usekomafont{sectioning}}
%\setkomafont{sectionentry}{\usekomafont{sectioning}}

%%% --- Titlepage ---
%\setkomafont{subject}{}
%\setkomafont{subtitle}{}
%\setkomafont{title}{}



\addtokomafont{sectioning}{\color{sectioncolor}} % Farbe der Ueberschriften
\IfChapterDefined{%
	\addtokomafont{chapter}{\color{sectioncolor}} % Farbe der Ueberschriften
}
\renewcommand*{\raggedsection}{\raggedright} % Titelzeile linksbuendig, haengend
%
%% UeberSchriften (Chapter und Sections) =================================

%%% Remove Space above Chapter.
%%% (NOT recommanded!)
%% Space above Chapter Title
% \renewcommand*{\chapterheadstartvskip}{\vspace{1\baselineskip}}%
%% Space below Chapter Title
% \renewcommand*{\chapterheadendvskip}{\vspace{0.5\baselineskip}}%

% -- Ueberschriften komlett Umdefinieren --
%%% Doc: ftp://tug.ctan.org/pub/tex-archive/macros/latex/contrib/titlesec/titlesec.pdf
%\usepackage{titlesec}

% -- Section Aussehen veraendern --
% --------------------------------
%% -> Section mit Unterstrich
% \titleformat{\section}
%   [hang]%[frame]display
%   {\usekomafont{sectioning}\Large}
%  {\thesection}
%   {6pt}
%   {}
%   [\titlerule \vspace{0.5\baselineskip}]
% --------------------------------

% -- Chapter Aussehen veraendern --
% --------------------------------
%--> Box mit (Kapitel + Nummer ) +  Name
% \titleformat{\chapter}[display]     % {command}[shape]
%   {\usekomafont{chapter}\filcenter} % format
%   {                                 % label
%   {\fcolorbox{black}{shadecolor}{
%   {\huge\chaptertitlename\mbox{\hspace{1mm}}\thechapter}
%   }}}
%   {1pc}                             % sep (from chapternumber)
%   {\vspace{1pc}}                    % {before}[after] (before chaptertitle and after)
% --------------------------------
%--> Kapitel + Nummer + Trennlinie + Name + Trennlinie
%\titleformat{\chapter}[display]	% {command}[shape]
%  {\usekomafont{chapter}\Large \color{black}}	% format
%  {   										% label
%  \LARGE\MakeUppercase{\chaptertitlename} \Huge \thechapter \filright%
%  }%}
%  {1pt}										% sep (from chapternumber)
%  {\titlerule \vspace{0.9pc} \filright \color{sectioncolor}}   % {before}[after] (before chaptertitle and after)
%  [\color{black} \vspace{0.9pc} \filright {\titlerule}]


%%% Doc: No documentation
% Indent first paragraph after section header
% \usepackage{indentfirst}

%% Captions (Schrift, Aussehen) ==========================================

% % Folgende Befehle werden durch das Paket caption und subfig ersetzt !
% \setcapindent{1em} % Einrueckung der Beschriftung
% \setkomafont{caption}{\color{black}\small\sffamily\RaggedRight}  % Schrift fuer Caption
% \setkomafont{captionlabel}{\color{black}\small}   % Schrift fuer 'Abbildung' usw.

%%% Doc: ftp://tug.ctan.org/pub/tex-archive/macros/latex/contrib/caption/caption.pdf
\usepackage{caption}
% Aussehen der Captions
\captionsetup{
   margin = 10pt,
   font = {small,rm},
   labelfont = {small,bf},
   format = plain, % oder 'hang'
   indention = 0em,  % Einruecken der Beschriftung
   labelsep = colon, %period, space, quad, newline
   justification = RaggedRight, % justified, centering
   singlelinecheck = true, % false (true=bei einer Zeile immer zentrieren)
   position = bottom %top
}
%%% Bugfix Workaround
\DeclareCaptionOption{parskip}[]{}
\DeclareCaptionOption{parindent}[]{}

% Aussehen der Captions fuer subfigures (subfig-Paket)
\IfPackageLoaded{subfig}{
 \captionsetup[subfloat]{%
   margin = 10pt,
   font = {small,rm},
   labelfont = {small,bf},
   format = plain, % oder 'hang'
   indention = 0em,  % Einruecken der Beschriftung
   labelsep = space, %period, space, quad, newline
   justification = RaggedRight, % justified, centering
   singlelinecheck = true, % false (true=bei einer Zeile immer zentrieren)
   position = bottom, %top
   labelformat = parens % simple, empty % Wie die Bezeichnung gesetzt wird
 }
}

% Aendern der Bezeichnung fuer Abbildung und Tabelle
% \addto\captionsngerman{% "captionsgerman" fuer alte  Rechschreibung
%   \renewcommand{\figurename}{Abb.}%
%   \renewcommand{\tablename}{Tab.}%
% }

% Caption fuer nicht fliessende Umgebungen
%%% Doc: ftp://tug.ctan.org/pub/tex-archive/macros/latex/contrib/misc/capt-of.sty
\IfPackageNotLoaded{caption}{
	\usepackage{capt-of} % only load when caption is not loaded. Otherwise compiling will fail.
	%Usage: \captionof{table}[short Titel]{long Titel}
}
%


%%% Doc: ftp://tug.ctan.org/pub/tex-archive/macros/latex/contrib/mcaption/mcaption.pdf
% Captions in Margins
% \usepackage[
% 	top,
% 	bottom
% ]{mcaption}

%%% Example:
% \begin{figure}
%   \begin{margincap}[short caption]{margin caption}
%     \centering
%     \includegraphics{picture}
%   \end{margincap}
% \end{figure}



% \numberwithin{figure}{chapter} %Befehl zum Kapitelweise Nummerieren der Bilder, setzt `amsmath' vorraus
% \numberwithin{table}{chapter}  %Befehl zum Kapitelweise Nummerieren der Tabellen, setzt `amsmath' vorraus

%% Inhaltsverzeichnis (Schrift, Aussehen) sowie weitere Verzeichnisse ====

\setcounter{secnumdepth}{2}    % Abbildungsnummerierung mit groesserer Tiefe
\setcounter{tocdepth}{2}		 % Inhaltsverzeichnis mit groesserer Tiefe
%

% Inhalte von List of Figures
\IfPackageLoaded{subfig}{
	\setcounter{lofdepth}{1}  %1 = nur figures, 2 = figures + subfigures
}

% -------------------------------------------------------

% Aussehen des Inhaltsverzeichnisses: tocloft
%%% Doc: ftp://tug.ctan.org/pub/tex-archive/macros/latex/contrib/tocloft/tocloft.pdf
%% Laden mit Option subfigure in Abhaengigkeit vom Paket subfigure und subfig
% \IfElsePackageLoaded{subfig}
% 	% IF subfig
% 	{\usepackage[subfigure]{tocloft}}{
% 	% ELSE
% 	\IfElsePackageLoaded{subfigure}
% 		% IF subfigure
% 		{\usepackage[subfigure]{tocloft}}
% 	   % Else (No subfig nor subfigure)
% 		{\usepackage{tocloft}}
% 	}
%
% %TOCLOFT zerstoert Layout der Ueberschriften von TOC, LOT, LOF
% \IfPackageLoaded{tocloft}{
% %
% %%%% Layout Matthias Pospiech (alles serifenlos)
% \IfChapterDefined{%
% 	\renewcommand{\cftchappagefont}{\bfseries\sffamily}  % Kapitel Seiten Schrift
% 	\renewcommand{\cftchapfont}{\bfseries\sffamily}      % Kapitel Schrift
% }
% \renewcommand{\cftsecpagefont}{\sffamily}            % Section Seiten Schrift
% \renewcommand{\cftsubsecpagefont}{\sffamily}         % Subsectin Seiten Schrift
% \renewcommand{\cftsecfont}{\sffamily}                % Section Schrift
% \renewcommand{\cftsubsecfont}{\sffamily}             % Subsection Schrift
%
% %%%% Layout aus Typokurz:
% % % Seitenzahlen direkt hinter TOC-Eintrag:
% % % Ebene \chapter
% % \renewcommand{\cftchapleader}{}
% % \renewcommand{\cftchapafterpnum}{\cftparfillskip}
% % % Ebene \section
% % \renewcommand{\cftsecleader}{}
% % \renewcommand{\cftsecafterpnum}{\cftparfillskip}
% % % Ebene \subsection
% % \renewcommand{\cftsubsecleader}{}
% % \renewcommand{\cftsubsecafterpnum}{\cftparfillskip}
% % % Abstaende vor Eintraegen im TOC verkleinern
% % \setlength{\cftbeforesecskip}{.4\baselineskip}
% % \setlength{\cftbeforesubsecskip}{.1\baselineskip}
% }
% % Ende tocloft Einstellungen --------------

%%% Doc: ftp://tug.ctan.org/pub/tex-archive/macros/latex/contrib/ms/multitoc.dvi
% TOC in mehreren Spalten setzen
%\usepackage[toc]{multitoc}

% -------------------------

%%Schriften fuer Minitoc (Inhaltsverzeichnis vor jedem Kapitel)
%%% Doc: ftp://tug.ctan.org/pub/tex-archive/macros/latex/contrib/minitoc/minitoc.pdf
%\IfElseChapterDefined{%
% \usepackage{minitoc}
% \setlength{\mtcindent}{0em} % default: 24pt
% \setcounter{minitocdepth}{2}
% \setlength{\mtcskipamount}{\bigskipamount}
% \mtcsettitlefont{minitoc}{\normalsize\SectionFontStyle}
% \mtcsetfont{minitoc}{*}{\small\SectionFontStyle} %\color{textcolor}
% \mtcsetfont{minitoc}{section}{\small\SectionFontStyle}
% \mtcsetfont{minitoc}{subsection}{\small\SectionFontStyle}
% \mtcsetfont{minitoc}{subsubsection}{\small\SectionFontStyle}
%}{
% \usepackage{minitoc}
% \setlength{\stcindent}{0pt} %default
% \setcounter{secttocdepth}{2} %default
% \mtcsettitlefont{secttoc}{\SectionFontStyle}
% \mtcsetfont{secttoc}{*}{\small\SectionFontStyle}%
% \mtcsetfont{secttoc}{subsection}{\small\SectionFontStyle}
% \mtcsetfont{secttoc}{subsubsection}{\small\SectionFontStyle}
%}

% Packages that MUST be loaded before minitoc !
% hyperref, caption, sectsty, varsects, fncychap, hangcaption, quotchap, romannum, sfheaders, alnumsec, captcont


%% Index & Co. ===========================================================
% gibts dafuer noch eine sauberere Loesung ?
%%%%%%%% Index zweispaltig %%%%%%%
% \makeatletter
% \renewenvironment{theindex}{%
% \setlength{\columnsep}{2em}
% \begin{multicols}{2}[\section*{\indexname}]
% \parindent\z@
% \parskip\z@ \@plus .3\p@\relax
% \let\item\@idxitem}%
% {\end{multicols}\clearpage}
% \makeatother
%%%%%%%%%%%%%%%%%%%%%%%%%%%%%%%%%%


% ~~~~~~~~~~~~~~~~~~~~~~~~~~~~~~~~~~~~~~~~~~~~~~~~~~~~~~~~~~~~~~~~~~~~~~~~
% PDF related packages
% ~~~~~~~~~~~~~~~~~~~~~~~~~~~~~~~~~~~~~~~~~~~~~~~~~~~~~~~~~~~~~~~~~~~~~~~~

%%% Doc: ftp://tug.ctan.org/pub/tex-archive/macros/latex/contrib/microtype/microtype.pdf
% Optischer Randausgleich mit pdfTeX
\ifpdf
\usepackage[%
	expansion=true, % better typography, but with much larger PDF file.
	protrusion=true
]{microtype}
\fi


%% Use only instead of hyperref !
% \usepackage[%
%    %ref,     % verweist auf Abschnitte
%    pageref, % verweist auf Seiten
% ]{backref} % Links in BiB back to Citation page/section (can be loaded by hyperref too)


%%% Doc: ftp://tug.ctan.org/pub/tex-archive/macros/latex/contrib/hyperref/doc/manual.pdf
\usepackage[
%   % Farben fuer die Links
%   colorlinks=true,         % Links erhalten Farben statt Kaeten
%   urlcolor=pdfurlcolor,    % \href{...}{...} external (URL)
%   filecolor=pdffilecolor,  % \href{...} local file
%   linkcolor=pdflinkcolor,  %\ref{...} and \pageref{...}
%   citecolor=pdfcitecolor,  %
%   % Links
%   raiselinks=true,			 % calculate real height of the link
%   breaklinks,              % Links berstehen Zeilenumbruch
%   backref=page,            % Backlinks im Literaturverzeichnis (section, slide, page, none)
%   pagebackref=true,        % Backlinks im Literaturverzeichnis mit Seitenangabe
%   verbose,
%   hyperindex=true,         % backlinkex index
   linktocpage=true,        % Inhaltsverzeichnis verlinkt Seiten
%   hyperfootnotes=false,     % Keine Links auf Fussnoten
%   % Bookmarks
   bookmarks=true,          % Erzeugung von Bookmarks fuer PDF-Viewer
%   bookmarksopenlevel=1,    % Gliederungstiefe der Bookmarks
   bookmarksopen=true,      % Expandierte Untermenues in Bookmarks
   bookmarksnumbered=true,  % Nummerierung der Bookmarks
%   bookmarkstype=toc,       % Art der Verzeichnisses
%   % Anchors
%   plainpages=false,        % Anchors even on plain pages ?
%   pageanchor=true,         % Pages are linkable
%   % PDF Informationen
   pdftitle={},             % Titel
   pdfauthor={Alexander Becker},            % Autor
%   pdfcreator={LaTeX, hyperref, KOMA-Script}, % Ersteller
%   %pdfproducer={pdfeTeX 1.10b-2.1} %Produzent
%   pdfdisplaydoctitle=true, % Dokumententitel statt Dateiname im Fenstertitel
%   pdfstartview=FitH,       % Dokument wird Fit Width geaefnet
%   pdfpagemode=UseOutlines, % Bookmarks im Viewer anzeigen
%   pdfpagelabels=true,           % set PDF page labels
%   pdfpagelayout=TwoPageRight, % zweiseitige Darstellung: ungerade Seiten
%   									 % rechts im PDF-Viewer
%   %pdfpagelayout=SinglePage, % einseitige Darstellung
]{hyperref}

\IfPackageLoaded{backref}{
   % % Change Layout of Backref
   \renewcommand*{\backref}[1]{%
   	% default interface
   	% #1: backref list
   	%
   	% We want to use the alternative interface,
   	% therefore the definition is empty here.
   }%
   \renewcommand*{\backrefalt}[4]{%
   	% alternative interface
   	% #1: number of distinct back references
   	% #2: backref list with distinct entries
   	% #3: number of back references including duplicates
   	% #4: backref list including duplicates
   	\mbox{(Zitiert auf %
   	\ifnum#1=1 %
		   Seite~%
	   \else
   		Seiten~%
   	\fi
   	#2)}%
   }
}

%%% Doc: ftp://tug.ctan.org/pub/tex-archive/macros/latex/contrib/oberdiek/hypcap.pdf
% Links auf Gleitumgebungen springen nicht zur Beschriftung,
% sondern zum Anfang der Gleitumgebung
\IfPackageLoaded{hyperref}{%
	\usepackage[figure,table]{hypcap}
}

% Auch Abbildung und nicht nur die Nummer wird zum Link (abgeleitet
% aus Posting von Heiko Oberdiek (d09n5p$9md$1@news.BelWue.DE);
% Verwendung: In \abbvref{label} ist ein Beispiel dargestellt
\providecommand*{\figrefname}{Abbildung }
\newcommand*{\figref}[1]{%
  \hyperref[fig:#1]{\figrefname{}}\ref{fig:#1}%
}
% ebenso bei Tabellen
\providecommand*{\tabrefname}{Tabelle~}
\newcommand*{\tabref}[1]{%
  \hyperref[tab:#1]{\tabrefname{}}\ref{tab:#1}%
}
% und Abschnitten
\providecommand*{\secrefname}{Abschnitt }
\newcommand*{\secref}[1]{%
  \hyperref[sec:#1]{\secrefname{}}\ref{sec:#1}%
}
% und Kapiteln
\providecommand*{\chaprefname}{Kapitel~}
\newcommand*{\chapref}[1]{%
  \hyperref[chap:#1]{\chaprefname{}}\ref{chap:#1}%
}

%%% Doc: ftp://tug.ctan.org/pub/tex-archive/macros/latex/contrib/pdfpages/pdfpages.pdf
\usepackage{pdfpages} % Include pages from external PDF documents in LaTeX documents

%%% Doc: ftp://tug.ctan.org/pub/tex-archive/macros/latex/contrib/oberdiek/pdflscape.sty
%\usepackage{pdflscape} %  Querformat mit PDF
%
% Pakete Laden die nach Hyperref geladen werden sollen
\LoadPackagesNow % (ltxtable, tabularx)

% ~~~~~~~~~~~~~~~~~~~~~~~~~~~~~~~~~~~~~~~~~~~~~~~~~~~~~~~~~~~~~~~~~~~~~~~~
% Zusätzliche Pakete
% ~~~~~~~~~~~~~~~~~~~~~~~~~~~~~~~~~~~~~~~~~~~~~~~~~~~~~~~~~~~~~~~~~~~~~~~~

% Todo Notes Package
% ----------------
% Weitere Informationen: http://www.tex.ac.uk/tex-archive/macros/latex/contrib/todonotes/todonotes.pdf
\usepackage{todonotes}

%%% Doc: ftp://tug.ctan.org/pub/tex-archive/macros/latex/contrib/hyphenat/hyphenat.pdf
% According to documentation the font warnings can be ignored
%\usepackage[htt]{hyphenat} % enable hyphenation of typewriter text word (\textt).

%% Komprimierung von Bildern in PDF ausschalten
% \ifpdf
%    \pdfcompresslevel=0
% \fi
% ~~~~~~~~~~~~~~~~~~~~~~~~~~~~~~~~~~~~~~~~~~~~~~~~~~~~~~~~~~~~~~~~~~~~~~~~
% end of preambel
% ~~~~~~~~~~~~~~~~~~~~~~~~~~~~~~~~~~~~~~~~~~~~~~~~~~~~~~~~~~~~~~~~~~~~~~~~
%\IfPackageLoaded{fancyvrb}{
%	\DefineShortVerb{\|} % Nur mit fancyvrb zusammen laden!
%}


% Auszufuehrende Befehle  ------------------------------------------------
\IfDefined{makeindex}{\makeindex}
\IfDefined{makenomenclature}{\makenomenclature}
\IfPackageLoaded{minitoc}{\IfElseUnDefined{chapter}{\dosecttoc}{\dominitoc}}
\renewcommand{\nomname}{Abkürzungsverzeichnis}

\listfiles
%------------------------------------------------------------------------

\setlength{\parindent}{0em}
\setlength{\parskip}{1em}

%
%%%% Neue Befehle
% % redefine \textmu to other mu commands usefull inside text
% \renewcommand{\textmu}{$\upmu$}

% Caption with defined width
\newcommand{\wcaption}[2]{%
   \begin{minipage}{#1}%
   \caption{#2}%
   \end{minipage}%
}

\newcommand{\figureref}[1]{(Abbildung \ref{#1})}%
\newcommand{\eqnref}[1]{(\ref{#1})}%

% Command for margin text with usefull style
%\newcommand{\marginlabel}[1]{\mbox{}\marginline{\hspace{0pt}\footnotesize\sffamily #1}}%
\newcommand{\marginlabel}[1]{\marginnote{#1}}%

%\newcommand{\comment}[1]{\marginnote{#1}}%

% Enable space for figures that extent into the margin (right and/or leftside)
% Can be used inside a figure
% Note: sidecap defines a similar environment 'wide' !
\newenvironment{widespace}[2]{%
   \begin{list}{}{%
      \setlength{\topsep}{0pt}%
      \setlength{\leftmargin}{#1}%
      \setlength{\rightmargin}{#2}%
      \setlength{\listparindent}{\parindent}%
      \setlength{\itemindent}{\parskip}%
   }%
   \item[]%
}%
{%
   \end{list}%
}%

\newlength{\marginwidth}
\setlength{\marginwidth}{\marginparwidth}
\addtolength{\marginwidth}{\marginparsep}

%% Beispiel:
% \begin{figure}
% \begin{widespace}{-\marginwidth}{0pt}
%  \subfloat[Bergzebrastute]
%  {\includegraphics[width=0.45\linewidth]{../Bilder/Eingewoehnung2.jpg}}
%  \hspace*{1em}
%  \subfloat[Morro Moco]
%  {\includegraphics[width=0.45\linewidth]{../Bilder/bergzebra2.jpg}}
% \end{widespace}
% \end{figure}


% quantum optics - Latex Commands: Math **********************************
% ------------------------------------------------------------------------
% by: Matthias Pospiech
%%%%%%%%%%%%%%%%%%%%%%%%%%%%%%%%%%%%%%%%%%%%%%%%%%%%%%%%%%%%%%%%%%%%%%%%%%


% --| Math |-------------------------------------------------------

% -- Replacements --
\newcommand{\comp}{\ast}
%\renewcommand{\dagger}{+}


% -- new commands --
\providecommand{\abs}[1]{\lvert#1\rvert}
\providecommand{\Abs}[1]{\left\lvert#1\right\rvert}
\providecommand{\norm}[1]{\left\Vert#1\right\Vert}
\providecommand{\Trace}[1]{\ensuremath{\Tr\{\,#1\,\}}} % Trace /Spur
%

% -- differentials --
\renewcommand{\d}{\partial\mspace{2mu}} % partial diff
\newcommand{\td}{\,\mathrm{d}}	% total diff
\newcommand{\ddt}[1]{\frac{\td #1}{\td t}}

% -- Abbrevitations --
\renewcommand{\Re}{\text{Re}}			% Real value
\renewcommand{\Im}{\text{Im}}			% Real value
\newcommand{\complex}{\mathbb{C}} % Complex
\newcommand{\real}{\mathbb{R}}    % Real
%\newcommand{\R}{\real}						% Real
%\newcommand{\N}{\mathbb{N}}
%\newcommand{\Z}{\mathbb{Z}}
\renewcommand{\L}{\mathcal{L}}
\newcommand{\N}{\mathcal{N}}
\newcommand{\R}{\mathcal{R}}
\newcommand{\D}{\mathcal{D}}
%
\newcommand{\Ham}{\mathcal{H}}    % Hamilton
\newcommand{\Prob}{\mathscr{P}}    % Hamilton
\newcommand{\unity}{\mathds{1}}   % Real
%

\newcommand\gammab{\gamma_\bot}
\newcommand\gammap{\gamma_\parallel}
\newcommand\gammai{\gamma_\text{int}}
\newcommand\gammae{\gamma_\text{ext}}

% -- New Operators --
\DeclareMathOperator{\rot}{rot}
\DeclareMathOperator{\grad}{grad}
%\DeclareMathOperator{\div}{div}
\renewcommand{\div}{\text{div}\,}
\DeclareMathOperator{\Tr}{Tr}
\DeclareMathOperator{\const}{const}
\DeclareMathOperator{\e}{e} 			% exponatial Function

% -- new symbols --
\newcommand{\laplace}{\Delta}
\newcommand{\dalembert}{\Box}

% -- new arrows --
\renewcommand{\leadsto}{\Longrightarrow}
\newcommand{\leftrightleadsto}{\Longleftrightarrow}


% -- Text subscripts--
\newcommand{\rel}{_\text{rel}}
%\newcommand{\st}{\text{st}}
%

% -- other --
\newcommand{\com}[2]{\underbrace{#1}_{\textrm{\scriptsize #2}}}
\newcommand{\with}[1]{\stackrel{\ref{#1}}{\Longrightarrow}}
%\newcommand{\unit}[1]{\,\textrm{#1}}

%\newcommand{\variance}[1]{\delta \mean{#1}^2}
\newcommand{\variance}[1]{(\Delta{#1})^2}
%\newcommand{\variance}[1]{\delta #1^2}

% -- Physics --------------------------------
\newcommand\op[1]{{\hat{\mathrm{#1}}}}  % Operator

\newcommand\expect[1]{\ensuremath{\left\langle{#1}\right\rangle}} %
%
\newcommand{\mean}[1]{\ensuremath{\overline{#1}}} % mean value
%
\newcommand{\state}[1]{\ensuremath{\ket{#1}}}
%
\newcommand\commutator[2]{\ensuremath{\mathinner{%
    \mathopen[\,#1,#2\,\mathclose]}}}
\newcommand{\Commutator}[2]{\ensuremath{\left[\,#1,#2\,\right]}}
\newcommand{\bigcommutator}[2]{\ensuremath{\bigl[\,#1,#2\,\bigr]}}
\newcommand{\Bigcommutator}[2]{\ensuremath{\Bigl[\,#1,#2\,\Bigr]}}
%
\newcommand\poisson[2]{\mathinner{%
    \mathopen\{#1,#2\mathclose\}}}
%

% -- Layout --------------------------------

\newcommand*{\dashfill}{\leavevmode\cleaders\hbox{-}\hfill\kern0pt}

\newcommand*{\midhrulefill}{
\leavevmode
\cleaders\hbox to 1ex{\raisebox{.5ex}{\rule{1ex}{.4pt}}}\hfill\kern0pt
}


%% Kommandos fuer Tabellen. Entnommen aus The LateX Companion, tabsatz.ps und diversen Dokus:

%%% ---| Farben fuer Tabellen |-------------------
\IfPackageLoaded{xcolor}{
   \colorlet{tablesubheadcolor}{gray!30}
   \colorlet{tableheadcolor}{gray!25}
   \colorlet{tableblackheadcolor}{black!100}
   \colorlet{tablerowcolor}{gray!10.0}
}
%%% ---------------------------------------------


%%% -| Neue Spaltendefinitionen 'columntypes' |--
%
% Belegte Spaltentypen:
% l - links
% c - zentriert
% r - rechts
% p,m,b  - oben, mittig, unten
% X - tabularx Auto-Spalte

% um Tabellenspalten mit Flattersatz zu setzen, muss \\ vor
% (z.B.) \raggedright geschuetzt werden:
\newcommand{\PreserveBackslash}[1]{\let\temp=\\#1\let\\=\temp}


% Spalten mit Flattersatz und definierte Breite:
% m{} -> mittig
% p{} -> oben
% b{} -> unten
%
% Linksbuendig:
\newcolumntype{v}[1]{>{\PreserveBackslash\RaggedRight\hspace{0pt}}p{#1}}
\newcolumntype{M}[1]{>{\PreserveBackslash\RaggedRight\hspace{0pt}}m{#1}}
% % Rechtsbuendig :
% \newcolumntype{R}[1]{>{\PreserveBackslash\RaggedLeft\hspace{0pt}}m{#1}}
% \newcolumntype{S}[1]{>{\PreserveBackslash\RaggedLeft\hspace{0pt}}p{#1}}
% % Zentriert :
% \newcolumntype{Z}[1]{>{\PreserveBackslash\Centering\hspace{0pt}}m{#1}}
% \newcolumntype{A}[1]{>{\PreserveBackslash\Centering\hspace{0pt}}p{#1}}

\newcolumntype{Y}{>{\PreserveBackslash\RaggedLeft\hspace{0pt}}X}
%%% Spalten fuer Mathematik
%
% serifenlose Matheschrift
\newcolumntype{s}[1]{%
	>{\DC@{.}{,}{#1}\mathsf\bgroup}l%
	<{\egroup\DV@end}%
}

% Tabellenspaltentyp fuer den Kopf: (Farbe + Ausrichtung)
\newcolumntype{H}[1]{>{\columncolor{tableheadcolor}}l}

% aequivalent aus typokurz (fett+grau+links)
% \newcolumntype{H}{>{\fontseries{b}\selectfont%
%     \columncolor[gray]{.8}[6pt][0pt]}l}
%%% --------------------------------------------


%%% ---|Listen in Tabellen |--------------------
\newcommand{\removeindentation}{%
	\leftmargini=\labelsep%
	\advance\leftmargini by \labelsep%
}
%
\makeatletter
\newcommand\tableitemize{
	\@minipagetrue%
	\removeindentation
}
\makeatother
%%% --------------------------------------------

%%% ---|Layout der Tabellen |-------------------

% Neue Umgebung fuer Tabellen:

\newenvironment{Tabelle}[2][c]{%
  \tablestylecommon
  \begin{longtable}[#1]{#2}
  }
  {\end{longtable}%
  \tablerestoresettings
}


% Groesse der Schrift in Tabellen
\newcommand{\tablefontsize}{ \footnotesize}
\newcommand{\tableheadfontsize}{\footnotesize}

% Layout der Tabelle: Ausrichtung, Schrift, Zeilenabstand
\newcommand\tablestylecommon{%
  \renewcommand{\arraystretch}{1.4} % Groessere Abstaende zwischen Zeilen
  \normalfont\normalsize            %
  \sffamily\tablefontsize           % Serifenlose und kleine Schrift
  \centering%                       % Tabelle zentrieren
}

\newcommand{\tablestyle}{
	\tablestylecommon
	%\tablealtcolored
}

% Ruecksetzten der Aenderungen
\newcommand\tablerestoresettings{%
  \renewcommand{\arraystretch}{1}% Abstaende wieder zuruecksetzen
  \normalsize\rmfamily % Schrift wieder zuruecksetzen
}

% Tabellenkopf: Serifenlos+fett+schraeg+Schriftfarbe
\newcommand\tablehead{%
  \tableheadfontsize%
  \sffamily\bfseries%
  %\slshape
  %\color{white}
}

\newcommand\tablesubheadfont{%
  \tableheadfontsize%
  \sffamily\bfseries%
  \slshape
  %\color{white}
}


\newcommand\tableheadcolor{%
	%\rowcolor{tablesubheadcolor}
	%\rowcolor{tableblackheadcolor}
	\rowcolor{tableheadcolor}%
}

\newcommand\tablesubheadcolor{%
	\rowcolor{tablesubheadcolor}
	%\rowcolor{tableblackheadcolor}
}


\newcommand{\tableend}{\arrayrulecolor{black}\hline}

% Tabellenkopf (1=Spaltentyp, 2=Text)
% \newcommand{\tablehead}[2]{
%   \multicolumn{1}{#1@{}}{%
%     \raisebox{.1mm}{% Ausrichtung der Beschriftung
%       #2%
%     }\rule{0pt}{4mm}}% unsichtbare Linie, die die Kopfzeile hoeher macht
% }


\newcommand{\tablesubhead}[2]{%
  \multicolumn{#1}{>{\columncolor{tablesubheadcolor}}l}{\tablesubheadfont #2}%
}

% Tabellenbody (=Inhalt)
\newcommand\tablebody{%
\tablefontsize\sffamily\upshape%
}

\newcommand\tableheadshaded{%
	\rowcolor{tableheadcolor}%
}
\newcommand\tablealtcolored{%
	\rowcolors{1}{tablerowcolor}{white!100}%
}
%%% --------------------------------------------

\newlength{\mylen}
\newlength{\adjusthspace}

\newenvironment{tabularc}[2]
{%
	\setlength\mylen{#2/(#1)-\tabcolsep*2-\arrayrulewidth*(#1+1)/(#1)}%
	%\setlength{\adjusthspace}{((#2-1)/2)*\linewidth}
	%\par\noindent
	%\hspace*{-\the\adjusthspace}
	\begin{tabular}%{#2}%
		{*{#1}{v{\the\mylen}}}%
}
{\end{tabular}\par}


%%% Silbentrennung
\hyphenation{}

%% Dokument Beginn %%%%%%%%%%%%%%%%%%%%%%%%%%%%%%%%%%%%%%%%%%%%%%%%%%%%%%%%

\begin{document}
% Deckblatt
% \subject{Diplomarbeit \\ Universität <einfügen>}
% \title{<Titel einfügen>}
% \author{Alexander Becker}
% \date{<Datum einfügen>}
% \maketitle

\begin{titlepage}

\begin{center}
\includegraphics[width=1\textwidth]{images/Logo_THM_MNI.jpg}
\par\end{center}


\noindent \begin{center}
\vspace{1cm}

\par\end{center}

\noindent \begin{center}
\textsf{\textbf{\huge Bachelorthesis}}\textsf{}\\
\textsf{}\\
\textsf{\Large Titel der Abschlussarbeit}
\par\end{center}{\Large \par}

\vspace{1cm}


\begin{center}
zur Erlangung des akademischen Grades
\par\end{center}

\begin{center}
{\textsf{\large  Bachelor of Science}}
\par\end{center}{\large \par}

\begin{center}
eingereicht im Fachbereich Mathematik, Naturwissenschaften und Informatik an der Technische Hochschule Mittelhessen{\large \vspace{1cm}
}
\par\end{center}{\large \par}

\begin{center}
von
\par\end{center}

\begin{center}
Dieudonne Timma Meyatchie
\par\end{center}

\begin{center}
\today
\par\end{center}

\vspace{1.5cm}


Referent: Prof. Dr. Peter Kneisel 

Korreferent: Herr Dennis Priefer 

\end{titlepage}


\thispagestyle{empty}
\chapter*{Erklärung der Selbstständigkeit}
\thispagestyle{empty}
Hiermit versichere ich, die vorliegende Arbeit selbstständig verfasst und keine anderen als die angegebenen Quellen und Hilfsmittel benutzt sowie die Zitate deutlich kenntlich gemacht zu haben.
\vspace{4\baselineskip}\\
Gießen, den \today \hfill Dieudonne Timma Meyatchie
\vspace{4\baselineskip}\\
\clearpage
\mbox{}\thispagestyle{empty}



% Nicht vergessen zu Entfernen wenn alles Fertig:
%\listoftodos

% Abstract - Kurzfassung der Arbeit ohne Wertung
% Wie schreibe ich das Abstract?
% Bei der Vorbereitung des Abstracts sollten Sie sich an folgenden grundlegenden Punkten orientieren:

%% Motivation des Textes:
%% worin liegt die Bedeutung der entsprechenden Forschung, warum sollte der längere Text gelesen werden?

%% Fragestellung:
%% welche Fragestellung(en) versucht der Text zu beantworten, was ist der Umfang der Forschung, was sind die zentralen Argumente und Behauptungen?

%% Methodologie:
%% welche Methoden/Zugänge nutzt der Autor/die Autorin, auf welche empirische Basis stützt sich der Text?

%% Ergebnisse:
%% zu welchen Ergebnissen kam die Forschung, was sind die zentralen Schlussfolgerungen des Textes?

%% Implikationen:
%% welche Schlussfolgerungen ergeben sich aus dem Text für die Forschung, was fügt der Text unserem Wissen über das Thema hinzu?

\begin{abstract}
Lorem ipsum dolor sit amet, consectetur adipisicing elit, sed do eiusmod tempor incididunt ut labore et dolore magna aliqua. Ut enim ad minim veniam, quis nostrud exercitation ullamco laboris nisi ut aliquip ex ea commodo consequat. Duis aute irure dolor in reprehenderit in voluptate velit esse cillum dolore eu fugiat nulla pariatur. Excepteur sint occaecat cupidatat non proident, sunt in culpa qui officia deserunt mollit anim id est laborum.

Et harum quidem rerum facilis est et expedita distinctio. Nam libero tempore, cum soluta nobis est eligendi optio cumque nihil impedit quo minus id quod maxime placeat facere possimus, omnis voluptas assumenda est, omnis dolor repellendus.

Sed ut perspiciatis unde omnis iste natus error sit voluptatem accusantium doloremque laudantium, totam rem aperiam, eaque ipsa quae ab illo inventore veritatis et quasi architecto beatae vitae dicta sunt explicabo. Nemo enim ipsam voluptatem quia voluptas sit aspernatur aut odit aut fugit, sed quia consequuntur magni dolores eos qui ratione voluptatem sequi nesciunt.

\end{abstract}
\cleardoublepage
\frontmatter
\cleardoublepage
% Inhaltsverzeichnis in den PDF-Links eintragen
\pdfbookmark[1]{Inhaltsverzeichnis}{toc}
\tableofcontents

% Hauptteil
\mainmatter
\chapter{Einführung}
\begin{quote}

\end{quote} 

\section{Motivation}

In dieser Arbeit wollen wir eine MDD-Infrastruktur für Joomla entwickeln, um die Migration  zu erleichtern, und unsere Erweiterungen schneller und besser zu implementieren.

\section{Problembeschreibung}
Die Nutzung von Content Management System vereinfachen zwar die Gestaltung von Webapplikation, aber die Migration  zu    neuen Versionen oder die Verschmelzung von mehreren Instanz stellen von verschiedenen Schwierigkeiten. Deswegen haben wir uns entschieden, eine MDD-Infrastruktur für Joomla zu entwickeln.   

\section{Fragestellung}
\begin{itemize}
\item Wie ist  Joomla in cJSL umgesetzt?
\item Wie kann man einen Forward-Generator zur Erstellung eines Joomla-Instanz\cite{xtext} mit Hilfe von cJSL entwickeln?
\item Wie kann man die  Daten aus dem cJSl-Instanz Model  mit einem Joomla Pakete zusammenbringen?
\item Wie kann man die generierte Joomla Instanz installieren?
\end{itemize}

\section{Ziele dieser Arbeit}

Die Ziele dieser Arbeit sind:

\begin{itemize}
\item Erstellung einer Abbildung von cJSL zu Joomla
\item Implementierung eines Forward-Generator
\item Lösungen zur Integration von Daten aus dem cJSL-Instanz Model zu einem Joomla-Pakete und zur Installation der so generierten Joomla-Instanz
\item Vorstellung von verschiedenen Einsatzmöglichkeiten
\end{itemize}

\section{Abgrenzung}
\begin{itemize}
\item Der Generator begrenzt sich auf dem CMS Joomla 2.5
\item Wir betrachten nur die Core-Erweiterungen
\item Generiert nur Code für die Datenhaltung und Konfiguration Daten der neuen Instanz.
\end{itemize}

\section{Vorgehensweise}
Zuerst werden wir mit Hilfe von Eclipse Modeling Framework eine MDD-Plattform aufstellen, um den Generator in Xtend zu implemetieren. Dann analysieren  wir eine Joomla Pakete und eine laufende System, um alle Spezifikation zu merken. Danach untersuchen wie diese Joomla-Spezifikationen in cJSL umgesetzt sind. Mit diesem Ergebnis werden wir unseren Forward-Generator implementieren. Mit Hilfe dieses Forward-Generators können wir verschiedene Migration Strategie für Joomla umsetzen.

\section{Struktur der Arbeit}

Die Arbeit besteht aus fünf Kapiteln: im ersten Kapitel "Technologieanalyse" erläuterte  ich  die verschiedenen Technologie und Konzepten wie zum Beispiel MDD-Infrastruktur und EMF, die für diese Arbeit nötig waren. Im zweiten Kapitel " Umsetzung von cJSL in Joomla " untersuche ich  cJSL und setze ich es in Joomla um. Im dritten Kapitel "Entwicklung eines cJSL-Forward-Generator" implementiere ich  auf Basis von cJSL eines Forward-Generator für Joomla. Im vierten Kapitel " Anwendungsmöglichkeit und Migrationsstrategie" stelle ich  verschiedene Anwendungsmöglichkeiten und Migrationsstrategie für eine oder mehreren Joomla Instanzen. Im fünften Kapitel "Fazit" mache ich eine Zusammenfassung der Arbeit.


% Alles was man wissen muss um das Konzept zu verstehen.

\chapter{Technologieanalyse}

\chapter{Umsetzung von cJSL in Joomla}

\chapter{Entwicklung eines cJSL-Forward-Generator}

\chapter{Zusammenfassung}

\section{Fazit}

\section{Auswertung}

\section{Weitere Ansätze}

\section{Nächste Schritte}

\section{Ausblick}


% Anhang (Bibliographie darf im deutschen nicht in den Anhang!)
\bibliography{myBib}
\clearpage
% Symbolverzeichnis (Begriffserklärung)
% Damit die Begriffe hinzugefügt werden muss in Texmaker unter Optionen - Konfiguren 
% Der Makeindex Befehl folgendermaßen geändert werden: "/usr/texbin/makeindex" %.nlo -s nomencl.ist -o %.nls
\IfDefined{printindex}{\printindex}
\IfDefined{printnomenclature}{\printnomenclature}
% 'Symbolverzeichnis' ins Inhaltsverzeichnis
\addcontentsline{toc}{chapter}{Abkürzungsverzeichnis}
% Abbildungs- und Tabellenverzeichnis
\listoffigures
\listoftables
\lstlistoflistings.

%\end{longtable}
% Anhang
\appendix
% 'Anhang' ins Inhaltsverzeichnis
%\phantomsection
%\addcontentsline{toc}{part}{Anhang}

\chapter{Anhang 1}
%\begin{table}[h] 
%\centering
\textbf{Menu} \linebreak
\textbf{Page} \linebreak
In der Untersuchung von MenuItems in Joomla stellen wir fest, dass es verschiedene Page Type existiert,die dieselbe   Elementen enthalten. Diese Elemente kommen zwar nicht in allen Items, aber so oft genug , damit wir sie eine extra Entität  Menuentry options in der Grammmatik widmen. Sie besteht aus selbst den Entitäten "Link type Options", "Page Display Options", "Metadata Options"  . Ihre Abbildungen sehen folgenden aus:
\linebreak
\linebreak
\begin{minipage}{0.5\textwidth}
\begin{tabular}{|p{0.6\textwidth}|p{0.6\textwidth} |}
\hline
\textbf{CJSL} 
&  
\textbf{JOOMLA} \\ \hline
\begin{itemize}
\item Link Title Attribute
\item Link CSS Style
\item Link Image
\item Add Menu Title
	\begin{enumerate}
	\item[|-]no| yes
	\end{enumerate}
\end{itemize}
&
\begin{itemize}
\item menu-anchor\_title
\item menu-anchor\_css
\item menu\_image
\item menu\_text
	\begin{enumerate}
		\item[|-] 0 |1
	\end{enumerate}
\end{itemize} \\ \hline
\end{tabular}
\captionof{table}{Link Type Options}
\end{minipage}

\begin{minipage}{0.6\textwidth}
\begin{tabular}{|p{0.6\textwidth}|p{0.6\textwidth} |}
\hline
\textbf{CJSL} 
&  
\textbf{JOOMLA} \\ \hline
\begin{itemize}
\item Browser Page Title
	\item Show Page Heading
	\begin{enumerate}
	\item[|-]no| yes
	\end{enumerate}
	\item Page Heading
	\item Page Class
\end{itemize}
&
\begin{itemize}
	\item page\_title
	\item show\_page\_heading
	\begin{enumerate}
	\item[|-] 0 |1
	\end{enumerate}
	\item page\_heading
	\item pageclass\_sfx
\end{itemize} \\ \hline
\end{tabular}
\captionof{table}{Page Display Options}
\end{minipage}


\begin{minipage}{0.6\textwidth}
\begin{tabular}{|p{0.6\textwidth}|p{0.6\textwidth} |}
\hline
\textbf{CJSL} 
&  
\textbf{JOOMLA} \\ \hline
\begin{itemize}
\item Meta Description
    \item Meta Keywords
    \item Robots
    \begin{enumerate}
	    \item[-] use global
	    \item[-] index, follow
	    \item[-] noindex, follow
	    \item[-]index, nofollow
	    \item[-] noindex,nofollow
    \end{enumerate}
    \item Secure
    \begin{enumerate}
	    \item[-]off| on| ignore
    \end{enumerate}
\end{itemize}
&
\begin{itemize}
\item menu-meta\_description
\item menu-meta\_keywords
\item robots
	\begin{enumerate}
	\item[-] Leer
	\item[-] index,follow
	\item[-] noindex, follow
	\item[-] index,nofollow
	\item[-] noindex,nofollow
	\end{enumerate}
\item secure
	\begin{enumerate}
	\item[-]-1| 1| 0
	\end{enumerate}
\end{itemize} \\ \hline
\end{tabular}
\captionof{table}{Metadata Options}
\end{minipage}
\linebreak
\\
\textbf{Articelpage}
\\
\\
In dem Contianer CJSL\_Articlepage enthält verschiene Entitäten, die im Joomla Menuitem Type darstellen:

\begin{itemize}
\item Archived Articles
\item Single Article
\item List All Categories
\item Category Blog
\item Category List
\item Featured Articles
\item Create Article
\end{itemize} 

Sie besitzen nicht nur die obengenannte Itemoptions gemeinsam, sondern auch verschiedene Entitäten, die nur bei Article Page vorkommen:

\begin{itemize}
\item Category Options
\item Categories Options
\item Blog Layout Options
\item Ordering and Pagination Options
\item Article Layout Options
\item Integrations Options
\item ArticleSettingtype
\item OptionSettingtype
\item MaxLevel Setting
\end{itemize}

Hier haben wir die Abbildung von den obengenannten Elemente: Die ArticleSettingtype -> ArticleStettingvalue, die ArticleOptiontype -> ArticleOptionvalue, und optionSettingtype -> optionSettingValue


\begin{minipage}{0.5\textwidth}
\begin{tabular}{|c|c|}
\hline
\textbf{CSL} & \textbf{Joomla} \\ 
\hline
 use global & leer \\ 
 \hline
use article setting & use-article\\ 
\hline
hide & 0 \\ 
\hline
show & 1 \\ 
\hline
\end{tabular}
\captionof{table}{ArticleSettingtype}
\end{minipage}
\begin{minipage}{0.5\textwidth}
\begin{tabular}{|c|c|}
\hline
\textbf{CSL} & \textbf{Joomla} \\ 
\hline
use global & leer \\ 
\hline
hide & 0 \\ 
\hline
show & 1 \\  
\hline
\end{tabular}
\captionof{table}{OptionSettingtype}
\end{minipage}

\begin{minipage}{0.4\textwidth}
\begin{tabular}{|c|c|}
\hline
\textbf{CSL} & \textbf{Joomla} \\ 
\hline
use global & leer \\ 
\hline
all & -1 \\ 
\hline
none & 0 \\ 
\hline
1-5 & 1-5 \\
\hline
\end{tabular}
\captionof{table}{MaxlevelSettingtype}
\end{minipage}
\begin{minipage}{0.6\textwidth}
\begin{tabular}{|p{0.6\textwidth} | p{0.5\textwidth}|}
\hline
\textbf{CSL} & \textbf{Joomla} \\ 
\hline
\begin{itemize}
	\item Show Feed Link
	\begin{enumerate}
		\item[|-]OptionSettingtype
	\end{enumerate}
	\item For each feed item show
	\begin{enumerate}
	 \item[|-] use global
	 \item[|-] intro text
	 \item[|-] full text
	\end{enumerate}
\end{itemize}
 & 
\begin{itemize}
	\item show\_feed\_link
	 \item[|-]OptionSettingvalue
	\item feed\_summary
	\begin{enumerate}
		 \item[|-]leer
		 \item[|-]0
		 \item[|-]1
		\end{enumerate}
	
\end{itemize}
\\
\hline
\end{tabular}
\captionof{table}{Integration Options}
\end{minipage}

\begin{minipage}{0.6\textwidth}
\begin{tabular}{|p{0.5\textwidth}|p{0.6\textwidth}|}
\hline
\textbf{CSL} & \textbf{Joomla} \\ 
\hline
\begin{itemize}
	\item Leading Articles 
	\item Intro Articles
	\item Columns
	\item  Links
	\item Multi Column Order
	\begin{enumerate}
	  	\item[-] use global | down | across
	\end{enumerate} 
\end{itemize}
 & 
\begin{itemize}
	\item num\_leading\_articles 
	\item num\_intro\_articles 
	\item num\_columns 
	\item num\_links 
	\item multi\_column\_order
	\begin{enumerate}
		\item[-] leer| 0| 1
	\end{enumerate} 
\end{itemize}
\\
\hline
\end{tabular}
\captionof{table}{Blog Layout Options}
\end{minipage}

\begin{minipage}{0.7\textwidth}
\begin{tabular}{| p{0.7\textwidth}|p{0.6\textwidth}|}
\hline
\textbf{CSL} & \textbf{Joomla} \\ 
\hline
\begin{itemize}
	\item Show Top Level Category Description
	\begin{enumerate}
			 \item[-] OptionSettingtype
	\end{enumerate} 
	\item Subcategory Levels 
		\begin{enumerate}
				 \item[-] MaxlevelSettingtype
		\end{enumerate}
	\item Empty Categories 
		\begin{enumerate}
			 \item[-] OptionSettingtype
		\end{enumerate}
	\item Subcategories Descriptions
	 \begin{enumerate}
	 		 \item[-] OptionSettingtype
	 \end{enumerate}
\end{itemize}
 & 
\begin{itemize}
	\item show\_base\_description
	\begin{enumerate}
			 \item[-] OptionSettingvalue
	\end{enumerate} 
	\item maxLevelcat 
		\begin{enumerate}
				 \item[-] MaxlevelSettingvalue
		\end{enumerate}
	\item show\_empty\_categories\_cat 
		\begin{enumerate}
			 \item[-] OptionSettingvalue
		\end{enumerate}
	\item show\_subcat\_desc\_cat
	 \begin{enumerate}
	 		 \item[-] OptionSettingvalue
	 \end{enumerate}
\end{itemize}
\\
\hline
\end{tabular}
\captionof{table}{Categories Options}
\end{minipage}


\begin{minipage}{0.7\textwidth}
\begin{tabular}{|p{0.6\textwidth} | p{0.6\textwidth}|}
\hline
\textbf{CSL} & \textbf{Joomla} \\ 
\hline
\begin{itemize}
	\item Category Title 
		\begin{enumerate}
			\item[-] OptionSettingtype
		\end{enumerate} 
	\item Category Description
	\begin{enumerate}
		\item[-] OptionSettingtype
	\end{enumerate}
	\item Category Image
	\begin{enumerate}
		\item[-] OptionSettingtype
	\end{enumerate} 
	\item Subcategory Levels
	\begin{enumerate}
		\item[-] OptionSettingtype
	\end{enumerate} 
	\item Empty Categories
	\begin{enumerate}
			\item[-] OptionSettingtype
	\end{enumerate}
	\item Subcategories Descriptions
	\begin{enumerate}
		\item[-] OptionSettingtype
	\end{enumerate}
\end{itemize}
 & 
\begin{itemize}
	\item show\_category\_title 
		\begin{enumerate}
			\item[-] OptionSettingvalue
		\end{enumerate} 
	\item show\_description
	\begin{enumerate}
		\item[-] OptionSettingvalue
	\end{enumerate}
	\item show\_description\_image
	\begin{enumerate}
		\item[-] OptionSettingvalue
	\end{enumerate} 
	\item maxLevel
	\begin{enumerate}
		\item[-] OptionSettingvalue
	\end{enumerate} 
	\item show\_empty\_categories
	\begin{enumerate}
			\item[-] OptionSettingvalue
	\end{enumerate} 
	\item show\_subcat\_desc
	\begin{enumerate}
		\item[-] OptionSettingvalue
	\end{enumerate} 
\end{itemize}
\\
\hline
\end{tabular}
\captionof{table}{Category Options}
\end{minipage}

\begin{minipage}{0.6\textwidth}
\begin{tabular}{|p{0.6\textwidth} | p{0.6\textwidth}|}
\hline
\textbf{CSL} & \textbf{Joomla} \\ 
\hline
\begin{itemize}
	\item Display Select
		\begin{enumerate}
			\item[-] OptionSettingtype
		\end{enumerate}
	\item Filter Field
		\begin{enumerate}
			\item[-] OptionSettingtype
		\end{enumerate} 
	\item Table Headings
		\begin{enumerate}
			\item[-] OptionSettingtype
		\end{enumerate}
	\item Show Date 
		\begin{enumerate}
			\item[-] OptionSettingtype
		\end{enumerate} 
	\item Date Format 
	\item Show Hits in List
		\begin{enumerate}
			\item[-] OptionSettingtype
		\end{enumerate} 
	\item Show Author in List
		\begin{enumerate}
			\item[-] OptionSettingtype
		\end{enumerate} 
\end{itemize}
 & 
\begin{itemize}
	\item show\_pagination\_limit
		\begin{enumerate}
			\item[-] OptionSettingvalue
		\end{enumerate}
	\item filter\_field
		\begin{enumerate}
			\item[-] OptionSettingvalue
		\end{enumerate} 
	\item show\_headings
		\begin{enumerate}
			\item[-] OptionSettingvalue
		\end{enumerate}
	\item list\_show\_date
		\begin{enumerate}
			\item[-] OptionSettingvalue
		\end{enumerate} 
	\item date\_format 
	\item list\_show\_hits
		\begin{enumerate}
			\item[-] OptionSettingvalue
		\end{enumerate} 
	\item list\_show\_author
		\begin{enumerate}
			\item[-] OptionSettingvalue
		\end{enumerate} 	
\end{itemize}
\\
\hline
\end{tabular}
\captionof{table}{List Layout Options}
\end{minipage}

\begin{minipage}{0.7\textwidth}
\begin{tabular}{|p{0.7\textwidth} | p{0.6\textwidth}|}
\hline
\textbf{CSL} & \textbf{Joomla} \\ 
\hline
\begin{itemize}
\item Category Order
	\begin{enumerate}
		\item[|-] Use Global | No order
		\item[|-] Title Alphabetial
		\item[|-] Title Reverse Alphabetical
		\item[|-] Category Manager Order
	\end{enumerate}
\item  Article Order
	\begin{enumerate}
		\item[|-] Use Global
		\item[|-] Featured Article Order
		\item[|-] Most Recent First
		\item[|-] Oldest first
		\item[|-] Title Alphabetical
		\item[|-] Title Reverse Alphabetical
		\item[|-] Author Alphabetical
		\item[|-] Author Alphabetical Reverse
		\item[|-] Most Hits 
		\item[|-] Leasts Hits
		\item[|-] Article Manager Order 
	\end{enumerate}
\item Date for Ordering
	 \begin{enumerate}
	 \item[|-]use global| created| modified| published
	 \end{enumerate}
\item Pagination
	 \begin{enumerate}
	 	 \item[|-]use global|hide |show|auto
	 \end{enumerate}
\item Pagination Results
	\begin{enumerate}
	\item[|-]use global| hide| show
	\end{enumerate}
\end{itemize}
& 
\begin{itemize}
\item orderby\_pri
	\begin{enumerate}
		\item[|-] leer* | none
		\item[|-] alpha 
		\item[|-]ralpha
		\item[|-] order
	\end{enumerate}
\item orderby\_sec
	\begin{enumerate}
		\item[|-] leer* 
		\item[|-] front
		\item[|-] rdate 
		\item[|-] date
		\item[|-] alpha 
		\item[|-] ralpha
		\item[|-] author
		\item[|-] rauthor
		\item[|-] hits  
		\item[|-] rhits
		\item[|-] order
	\end{enumerate}
\item order\_date
	\begin{enumerate}
	\item[|-] leer | created | modified | published 
	\end{enumerate}
\item show\_pagination
	\begin{enumerate}
	\item[|-] leer | 0 | 1 | 2
	\end{enumerate}
\item show\_pagination\_results
	\begin{enumerate}
	\item[|-] leer | 0 | 1
	\end{enumerate}
\end{itemize}
\\
\hline
\end{tabular}
\captionof{table}{Ordering and Pagination Options}
\end{minipage}

\begin{minipage}{0.7\textwidth}
\begin{tabular}{|p {0.7\textwidth} | p {0.7\textwidth} |}
\hline
\textbf{CSL} & \textbf{Joomla} \\ 
\hline
\begin{itemize}
\item Show Title
	\begin{enumerate}
	\item[|-]ArticleSettingtype
	\end{enumerate}
\item Link Titles*
\item Show Intro*
\item Show Category*
\item Link Category*
\item Show Parent Category*
\item Link parent category*
\item Show Author*
\item Link Author*
\item Show create Date*
\item Show Modify Date*
\item Show Publish Date*
\item Show Item Navigation*
\item Show Vote*
\item Show Readmore*
\item Show Readmore Title*
\item Show Icons*
\item Show Print Icon*
\item Show Email Icon*
\item Show Hits*
\item Show Unauthorised Links*
\end{itemize} 
&
\begin{itemize}
\item show\_title
	\begin{enumerate}
     \item[-]ArticleSettingvalue
	\end{enumerate}	
\item link\_titles*
\item show\_intro*
\item show\_category*
\item link\_category*
\item show\_parent\_category*
\item link\_parent\_category*
\item show\_author*
\item link\_author*
\item show\_create\_date*
\item show\_modify\_date*
\item show\_publish\_date*
\item show\_item\_navigation*
\item show\_vote*
\item show\_readmore*
\item show\_readmore\_title*
\item show\_icons*
\item show\_print\_icon*
\item show\_email\_icon*
\item show\_hits*
\item show\_noauth*
\end{itemize}
\\ \hline
\end{tabular}
\captionof{table}{Article Settings Options}
\end{minipage}

* bedeudet, dass der Element ArticleSettingtype B.z AticleSettingvalue als Parameter hat.

\begin{minipage}{0.7\textwidth}
\begin{longtable}{| p {0.7\textwidth} |p {0.7\textwidth}|}
\hline
\textbf{CJSL} 
&  
\textbf{JOOMLA} \\ \hline
 1- List all Categories
  \begin{itemize}
    \item Basic  Setting 
   		\begin{enumerate}
   			\item[-] Select a Category
   		\end{enumerate}
   	\item Categories Options
	   	\begin{enumerate}
	    \item[+] Top Level Category Description
	    \item[+] Articles in Category 
	    	  \begin{enumerate}
	    			 \item[-] OptionSettingtype
	    	\end{enumerate}
	   \end{enumerate}
   	\item Category Options
   	\begin{enumerate}
   	\item[+] Articles in Category
   		\begin{enumerate}
   			\item[|-] OptionSettingtype
   		\end{enumerate}
   	\item[+] No Articles Message
   		\begin{enumerate}
   			\item[-] OptionSettingtype
   		\end{enumerate} 
   	\end{enumerate}
   	\item Blog Layout Options
   		\begin{enumerate}
   	   			\item[+] Include Subcategories
   			         \begin{enumerate}
   			          \item[|-] MaxlevelSettingType
   			         \end{enumerate}
   	   		\end{enumerate}
   	\item Ordering and Pagination Options
   	\item List Layout Options
   	\begin{enumerate}
   	\item[+] Articles to List
   			\begin{enumerate}
   				\item[|-] 5 bis 100 | all
   			\end{enumerate}
   	\end{enumerate}
   	\item Artikel Layout Options
   	\item Integration Options
   	\item Menuentry Options
 \end{itemize}
&
\textbf{.}List all Categories
	 \begin{itemize}
       \item Required  Setting 
      		\begin{enumerate}
      			\item[-] featured\_categories
      		\end{enumerate}
      	\item Categories Options
      		\begin{enumerate}
      		    \item[+] categories\_description
      		    \item[+] show\_cat\_num\_articles\_cat 
      		    	  \begin{enumerate}
      		    			 \item[|-] OptionSettingvalue
      		    	\end{enumerate}
      		   \end{enumerate}
   	   	\item Category Options
   	   		\begin{enumerate}
   	   	   	\item[+] show\_cat\_num\_articles
   	   	   		\begin{enumerate}
   	   	   			\item[|-] OptionSettingtype
   	   	   		\end{enumerate} 
   	   	   	\item[+] show\_no\_articles
   	   	   		\begin{enumerate}
   	   	   			\item[|-] OptionSettingvalue
   	   	   		\end{enumerate}
   	   	   	\end{enumerate}
   	   	\item Blog Layout Options
   	   		\begin{enumerate}
   	   	   	   			\item[+] show\_subcategory\_content
   	   	   			         \begin{enumerate}
   	   	   			          \item[|-] MaxlevelSettingvalue
   	   	   			         \end{enumerate}
   	   	   	   		\end{enumerate}
   	   	\item Ordering and Pagination Options
   	   	\item List Layout Options
   	   	\begin{enumerate}
   	   	\item[+] display\_num
   	   	\begin{enumerate}
   	   				\item[|-] 5 bis 100 | 0
   	   			\end{enumerate}
   	   	\end{enumerate}
   	   	\item Artikel Layout Options
   	   	\item Integration Options
   	   	\item Menuitem Options  
     \end{itemize}   
 	\\ \hline
2- Archived Article
\begin{itemize}
 \item Basic  Setting 
   		\begin{enumerate}
   		\item[-] Article Order (A.12) 
   		\item[-] Date for Ordering (A.12) 
   		\item[-] Articles to List (A.11)
   		\item[-] Filter Field (A.11) 
   		\item[-] Intro text Limit
   		\end{enumerate}
 	\item Artikel Layout Options
	\item Menuentry Options
\end{itemize}
&
2- Archived Article
\begin{itemize}
	\item Required  Setting 
   		\begin{enumerate}
   		\item[-] orderby\_sec (A.12)
   		\item[-] order\_date (A.12)
   		\item[-] display\_num (A.11) 
   		\item[-] filter\_field (A.11) 
   		\item[-] introtext\_limit 
   		\end{enumerate}
 	\item Artikel Layout Options
	\item Menuitem Options
\end{itemize}
\\ \hline
\end{longtable}
\captionof{table}{Articlepages I}
\end{minipage}

\begin{minipage}{0.8\textwidth}
\begin{longtable}{| p {0.7\textwidth} |p {0.7\textwidth}|}
\hline
\textbf{CJSL} 
&  
\textbf{JOOMLA} \\ \hline
4- Featured Articles
\begin{itemize}
	\item Required  Setting 
  		\begin{enumerate}
  		\item[-] Select Categories
  		\item[-] ...
		\end{enumerate}
	\item Blog Layout Options
   	\item Ordering and Pagination Options
   	\item Artikel Layout Options
   	\item Integration Options
   	\item Menuentry Options  
\end{itemize}
&
4- Featured Articles
\begin{itemize}
	\item Required  Setting 
  		\begin{enumerate}
  		\item[-] featured\_categories
  		\item[-] layout\_type
		\end{enumerate}
	\item Blog Layout Options
   	\item Ordering and Pagination Options
   	\item Artikel Layout Options
   	\item Integration Options
   	\item Menuitem Options  
\end{itemize}
\\ \hline
5-Category Blog Article
\begin{itemize}
\item Basic Setting
	\begin{enumerate}
		\item[-] Select Categories
		\item[-] .... 
	 \end{enumerate}
\item Category Options
	\begin{enumerate}
		\item[+] Show Subcategories Text
			\begin{enumerate}
			\item[|-] OptionSettingsType
			\end{enumerate}
		\item[+] Page Subheading
	   	\item[+] Articles in Category
	   		\begin{enumerate}
	   			\item[|-] OptionSettingtype
	   		\end{enumerate}
	   	\item[+] No Articles Message
	   		\begin{enumerate}
	   			\item[|-] OptionSettingtype
	   		\end{enumerate}
	\end{enumerate}
\item Blog Layout Options
	\begin{enumerate}
		\item[+] Include Subcategories
			\begin{enumerate}
			\item[|-] MaxlevelSettingType
			\end{enumerate}
	\end{enumerate}
\item Ordering and Pagination Options
\item Artikel Layout Options
\item Integration Options
\item Menuentry Options  
\end{itemize}
&
5- Category Blog Article
\begin{itemize}
\item Basic Setting
	\begin{enumerate}
		\item[-] featured\_categories
		\item[-] layout\_type 
	 \end{enumerate}
\item Category Options
	\begin{enumerate}
		\item[+] show\_category\_heading\_title\_text
			\begin{enumerate}
			\item[|-] OptionSettingsValue
			\end{enumerate}
		\item[+] page\_subheading
		\item[+] show\_cat\_num\_articles
	   		\begin{enumerate}
	   			\item[|-] OptionSettingtype
	   		\end{enumerate}
	   	\item[+] show\_no\_articles
	   		\begin{enumerate}
	   			\item[|-] OptionSettingvalue
	   		\end{enumerate} 
	\end{enumerate}
\item Blog Layout Options
	\begin{enumerate}
		\item[+] show\_subcategory\_content
			\begin{enumerate}
			\item[|-] MaxlevelSettingValue
			\end{enumerate}
	\end{enumerate}
\item Ordering and Pagination Options
\item Artikel Layout Options
\item Integration Options
\item Menuitem Options 
\end{itemize}
\\ \hline
\end{longtable}
\captionof{table}{Articlepages II}
\end{minipage}

\begin{minipage}{0.7\textwidth}
\begin{longtable}{| p {0.7\textwidth} |p {0.8\textwidth}|}
\hline
\textbf{CJSL} 
&  
\textbf{JOOMLA} \\ \hline
6-Category List Articles
\begin{itemize}
\item Basic Setting
	\begin{enumerate}
			\item[-] Featured Categories
	 \end{enumerate}
\item Category Options
	\begin{enumerate}
		\item[+]Show Category heading Title  
			\begin{enumerate}
			\item[|-] OptionSettingsType
			\end{enumerate}
		\item[+] Page Subheading
		\item[+] Articles in Category
	   		\begin{enumerate}
	   			\item[|-] OptionSettingtype
	   		\end{enumerate}
	   	\item[+] No Articles Message
	   		\begin{enumerate}
	   			\item[|-] OptionSettingtype
	   		\end{enumerate}  
	\end{enumerate}
\item List Layout Options
\begin{enumerate}
   	   	\item[+] Articles to List
   	   	\begin{enumerate}
   	   				\item[|-] 5 bis 100 | 0
   	   			\end{enumerate}
   	   	\end{enumerate}
\item Ordering and Pagination Options
\item Artikel Layout Options
\item Integration Options
\item Menuentry Options
\end{itemize}
&
6-Category List Articles
\begin{itemize}
\item Basic Setting
	\begin{enumerate}
			\item[-] featured\_categories
	 \end{enumerate}
\item Category Options
	\begin{enumerate}
		\item[+] show\_category\_heading\_title\_text
			\begin{enumerate}
			\item[|-] OptionSettingsValue
			\end{enumerate}
	 \item[+] page\_subheading
	 \item[+] show\_cat\_num\_articles
    		\begin{enumerate}
    			\item[|-] OptionSettingtype
    		\end{enumerate} 
	 \item[+] show\_no\_articles
	 	\begin{enumerate}
	 		\item[-] OptionSettingvalue
	 	\end{enumerate}  
	\end{enumerate}
\item List Layout Options
\begin{enumerate}
   	   	\item[+] display\_num
   	   	\begin{enumerate}
   	   				\item[|-] 5 bis 100 | 0
   	   			\end{enumerate}
   	   	\end{enumerate}
\item Ordering and Pagination Options
\item Artikel Layout Options
\item Integration Options
\item Menuitem Options 
\end{itemize}
\\ \hline
7-Create Article
\begin{itemize}
\item Basic Setting
	\begin{enumerate}
			\item[-] Default category 
			\item[-] Choose a category 
	 \end{enumerate}
\item Menuentry Options
\end{itemize}
&
7-Create Article
\begin{itemize}
\item Basic Setting
	\begin{enumerate}
			\item[-] enable\_category 
			\item[-] catid 
	 \end{enumerate}
\item Menuitem Options
\end{itemize}
\\ \hline
3-Single Article
\begin{itemize}
	\item Basic  Setting
		\begin{enumerate}
		\item[-] Select Article
		\end{enumerate}
	\item Artikel Layout Options
   	\item Integration Options
   	\item Menuentry Options
  \end{itemize}
&
3-Single Article
 \begin{itemize}
 	\item Required  Setting
 		\begin{enumerate}
 		\item[-] id
 		\end{enumerate}
 	\item Artikel Layout Options
    \item Integration Options
    \item Menuitem Options 	
\end{itemize}
\\ \hline
\end{longtable}
\captionof{table}{Articlepages III}
\end{minipage}



\textbf{UserPage}

Der Container USerpage besteht aus verschiedene Entitäten, die sich um die Darstellung des User widmen:
\begin{itemize}
	\item Login Form
	\item User Profile
	\item Edit User Profile
	\item Registration Form
	\item Username Reminder Request
	\item Password Reset
\end{itemize}

\begin{minipage}{0.7\textwidth}
\begin{longtable}{| p {0.7\textwidth} |p {0.8\textwidth}|}
\hline
\textbf{CJSL}
&  
\textbf{JOOMLA} \\ \hline
1- Login Form 
\begin{itemize}
\item Basic Options
	\begin{enumerate}
	\item[-]  Login Redirect
	\item[-]  Login Description
	\begin{enumerate}
	\item[|-]hide |show 
	\end{enumerate}
	\item[-]  Login Description Text 
	\item[-]  Login Image 
	\item[-]  Logout Redirect 
	\item[-]  Logout Text
	\begin{enumerate}
		\item[|-]hide | show
	\end{enumerate} 
	\item[-]  Logout Description Text 
	\item[-]  Logout image 
	\end{enumerate}
\item Menuentry Options
\end{itemize} 
&
1- Login Form 
\begin{itemize}
\item Basic Options
	\begin{enumerate}
	\item[-]  login\_redirect\_url
	\item[-]  logindescription\_show
	\begin{enumerate}
	\item[|-]0 | 1
	\end{enumerate}
	\item[-]  login\_description 
	\item[-]  login\_image 
	\item[-]  logout\_redirect\_url 
	\item[-]  logoutdescription\_show
	\begin{enumerate}
		\item[|-]0 | 1
	\end{enumerate} 
	\item[-]  logout\_description 
	\item[-]  logout\_image 
	\end{enumerate}
\item Menuitem Options
\end{itemize} 
\\ \hline
2- Login Form
\begin{itemize}
\item Menuentry Options
\end{itemize}
&
2- Login Form
\begin{itemize}
\item Menuitem Options
\end{itemize}
\\ \hline
\end{longtable}
\captionof{table}{Userpage I}
\end{minipage}

\begin{minipage}{0.7\textwidth}
\begin{longtable}{| p {0.7\textwidth} |p {0.8\textwidth}|}
\hline
\textbf{CJSL}
&  
\textbf{JOOMLA} \\ \hline
3- User Profile
\begin{itemize}
\item Menuentry Options
\end{itemize}
&
3-User Profile
\begin{itemize}
\item Menuitem Options
\end{itemize}
\\ \hline
4- Edit User Profile
\begin{itemize}
\item Menuentry Options
\end{itemize}
&
4-Edit User Profile
\begin{itemize}
\item Menuitem Options
\end{itemize}
\\ \hline
5-Registration Form
\begin{itemize}
\item Menuentry Options
\end{itemize}
&
5-Registration Form
\begin{itemize}
\item Menuitem Options
\end{itemize}
\\ \hline
6-Username Reminder Request
\begin{itemize}
\item Menuentry Options
\end{itemize}
&
6-Username Reminder Request
\begin{itemize}
\item Menuitem Options
\end{itemize}
\\ \hline
7-Password Reset
\begin{itemize}
\item Menuentry Options
\end{itemize}
&
7-Password Reset
\begin{itemize}
\item Menuitem Options
\end{itemize}
\\ \hline
\end{longtable}
\captionof{table}{Userpages II}
\end{minipage}



\textbf{Contacts}

Wie schon bei den Articlepage, stellen wir schnell fest, dass es verschiedenen wiederkehrende Elemente gibt, die in ENtitäten umwandeln kann :
\begin{itemize}
\item Categories Options
\item Category Options
\item List Layout Options
\item Contact Display Options
\item Form for Email Options
\end{itemize}

\begin{minipage}{0.7\textwidth}
\begin{longtable}{| p {0.7\textwidth} |p {0.8\textwidth}|}
\hline
\textbf{CJSL}
&  
\textbf{JOOMLA} \\ \hline
 \begin{enumerate}
	\item[-] Show Contact Form
		\begin{enumerate}
	   			\item[|-] OptionSettingtype
	   		\end{enumerate} 
	\item[-] Send Copy to Submitter
		\begin{enumerate}
	   			\item[|-] OptionSettingtype
	   		\end{enumerate}
	\item[-] Banned Email 
	\item[-] Banned Subject 
	\item[-] Banned Text
	\item[-] Session Check
		\begin{enumerate}
	   			\item[|-] OptionSettingtype
	   		\end{enumerate} 
	\item[-] Custom Reply
		\begin{enumerate}
	   			\item[|-] OptionSettingtype
	   		\end{enumerate} 
	\item[-] Contact Redirect
	\end{enumerate}
 &
 \begin{enumerate}
	\item[-] show\_email\_form
		\begin{enumerate}
	   			\item[|-] OptionSettingvalue
	   		\end{enumerate} 
	\item[-] show\_email\_copy
		\begin{enumerate}
	   			\item[|-] OptionSettingvalue
	   		\end{enumerate} 
	\item[-] banned\_email 
	\item[-] banned\_subject 
	\item[-] banned\_text 
	\item[-] validate\_session
		\begin{enumerate}
	   			\item[|-] OptionSettingtype
	   		\end{enumerate} 
	\item[-] custom\_reply
		\begin{enumerate}
	   			\item[|-] OptionSettingtype
	   		\end{enumerate} 
	\item[-] redirect
	\end{enumerate}
\\ \hline
\end{longtable}
\captionof{table}{Form for Email Options}
\end{minipage}


\begin{minipage}{0.6\textwidth}
\begin{tabular}{|p{0.6\textwidth} | p{0.6\textwidth}|}
\hline
\textbf{CSL} & \textbf{Joomla} \\ 
\hline
\begin{itemize}
		\item[-] Display Select
		\begin{enumerate}
	   			\item[|-] OptionSettingtype
	   		\end{enumerate}  
		\item[-] Table Headings
		\begin{enumerate}
   			\item[|-] OptionSettingtype
   		\end{enumerate} 
		\item[-] Position
			\begin{enumerate}
		   			\item[|-] OptionSettingtype
		   		\end{enumerate} 
		\item[-] Email
			\begin{enumerate}
		   			\item[|-] OptionSettingtype
		   		\end{enumerate} 
		\item[-] Phone
			\begin{enumerate}
		   			\item[|-] OptionSettingtype
		   		\end{enumerate} 
		\item[-] Mobile
			\begin{enumerate}
		   			\item[|-] OptionSettingtype
		   		\end{enumerate} 
		\item[-] Fax
			\begin{enumerate}
		   			\item[|-] OptionSettingtype
		   		\end{enumerate} 
		\item[-] City or Suburb
			\begin{enumerate}
		   			\item[|-] OptionSettingtype
		   		\end{enumerate} 
		\item[-] State or County
			\begin{enumerate}
		   			\item[|-] OptionSettingtype
		   		\end{enumerate} 
		\item[-] Country
			\begin{enumerate}
		   			\item[|-] OptionSettingtype
		   		\end{enumerate} 
		\item[-] Pagination
			\begin{enumerate}
		   			\item[|-] OptionSettingtype
		   		\end{enumerate} 
		\item[-] Pagination Results
			\begin{enumerate}
		   			\item[|-] OptionSettingtype
		   		\end{enumerate}
	\end{itemize}
 & 
\begin{itemize}
		\item[-] show\_pagination\_limit
		\begin{enumerate}
	   			\item[|-] OptionSettingvalue
	   		\end{enumerate}  
		\item[-] show\_headings
		\begin{enumerate}
		   			\item[|-] OptionSettingvalue
		   		\end{enumerate} 
		\item[-] show\_position\_headings
		\begin{enumerate}
		   			\item[|-] OptionSettingvalue
		   		\end{enumerate} 
		\item[-] show\_email\_headings
			\begin{enumerate}
			   			\item[|-] OptionSettingvalue
			   		\end{enumerate} 
		\item[-] show\_telephone\_headings
			\begin{enumerate}
			   			\item[|-] OptionSettingvalue
			   		\end{enumerate} 
		\item[-] show\_mobile\_headings
			\begin{enumerate}
			   			\item[|-] OptionSettingvalue
			   		\end{enumerate} 
		\item[-] show\_fax\_headings
			\begin{enumerate}
			   			\item[|-] OptionSettingvalue
			   		\end{enumerate} 
		\item[-] show\_suburb\_headings
			\begin{enumerate}
			   			\item[|-] OptionSettingvalue
			   		\end{enumerate} 
		\item[-] show\_state\_headings
			\begin{enumerate}
			   			\item[|-] OptionSettingvalue
			   		\end{enumerate} 
		\item[-] show\_country\_headings
			\begin{enumerate}
			   			\item[|-] OptionSettingvalue
			   		\end{enumerate} 
		\item[-] show\_pagination
			\begin{enumerate}
			   			\item[|-] OptionSettingvalue
			   		\end{enumerate} 
		\item[-] show\_pagination\_results
			\begin{enumerate}
			   			\item[|-] OptionSettingvalue
			   		\end{enumerate}
	\end{itemize}
\\
\hline
\end{tabular}
\captionof{table}{List Layout Options}
\end{minipage}

\begin{minipage}{0.6\textwidth}
\begin{longtable}{|p{0.6\textwidth} | p{0.6\textwidth}|}
\hline
\textbf{CSL} & \textbf{Joomla} \\ 
\hline
 \begin{itemize}
\item Display format
	\begin{enumerate}
	\item[|-] use Global | sliders
	\item[|-] tabs | plain
	\end{enumerate}
\item Contact Category
	\begin{enumerate}
	\item[|-] use Global | hide
	\item[|-] show without link | show with link
	\end{enumerate} 
\item show\_contact\_list
	\begin{enumerate}
		\item[|-] OptionSettingtype
	\end{enumerate} 
\item Name* 
\item Contact's Position* 
\item Email* 
\item Street Address* 
\end{itemize}
 & 
 \begin{itemize}
\item presentation\_style
\begin{enumerate}
	\item[|-] leer | sliders
	\item[|-] tabs | plain
	\end{enumerate} 
\item show\_contact\_category
\begin{enumerate}
	\item[|-]leer| hide
	\item[|-]show\_without\_link | show\_with\_link
	\end{enumerate} 
\item show\_contact\_list
\begin{enumerate}
		\item[|-] OptionSettingtvalue
	\end{enumerate} 
\item show\_name* 
\item show\_position* 
\item show\_email* 
\item show\_street\_address* 
\end{itemize}
\\
\hline
\end{longtable}
\end{minipage}

\begin{minipage}{0.6\textwidth}
\begin{tabular}{|p{0.6\textwidth} | p{0.6\textwidth}|}
\hline
\begin{itemize}
\item City or Suburb*
\item State or County* 
\item Postal Code* 
\item Country*
\item Telephone* 
\item Mobile phone* 
\item Fax*
\item Webpage* 
\item Misc. Information* 
\item Image* 
\item vCard* 
\item Show User Articles* 
\item Show Links* 
\item Link A Label 
\item Link B Label 
\item Link C Label 
\item Link D Label 
\item Link E Label
\end{itemize}
 & 
\begin{itemize}
\item show\_suburb* 
\item show\_state* 
\item show\_postcode* 
\item show\_country* 
\item show\_telephone* 
\item show\_mobile* 
\item show\_fax* 
\item show\_webpage* 
\item show\_misc* 
\item show\_image* 
\item allow\_vcard* 
\item show\_articles* 
\item show\_links* 
\item linka\_name 
\item linkb\_name 
\item linkc\_name 
\item linkd\_name 
\item linke\_name
\end{itemize}
\\ \hline
\end{tabular}
\captionof{table}{Contact Display Options}
\end{minipage}

\begin{minipage}{0.7\textwidth}
\begin{tabular}{|p{0.6\textwidth} | p{0.6\textwidth}|}
\hline
\textbf{CSL} & \textbf{Joomla} \\ 
\hline
 1- List all Categories
   \begin{itemize}
     \item Basic  Setting 
    		\begin{enumerate}
    			\item[-] Select a Category
    		\end{enumerate}
    	\item Categories Options
 	   	\begin{enumerate}
 	    \item[+] Top Level Category Description
 	    \item[+] Contacts in Category 
 	    	  \begin{enumerate}
 	    			 \item[-] OptionSettingtype
 	    	\end{enumerate}
 	   \end{enumerate}
    	\item Category Options
    	\item[+] Contacts in Category 
    		    	  \begin{enumerate}
    		    			 \item[-] OptionSettingtype
    		    	\end{enumerate}
    	\item Contact List Layouts
 	\item Contact Display Options
 	\item Form for Email Options
 	\item Integration Options
 	\item Menuentry Options
  \end{itemize}
 & 
1- List all Categories
  \begin{itemize}
    \item Basic  Setting 
   		\begin{enumerate}
   			\item[-] categorie\_id
   		\end{enumerate}
   	\item Categories Options
	   	\begin{enumerate}
	    \item[+] categories\_description
	    \item[+] show\_cat\_items\_cat 
	    	  \begin{enumerate}
	    			 \item[-] OptionSettingtype
	    	\end{enumerate}
	   \end{enumerate}
   	\item Category Options
   	\item[+] show\_cat\_items
   		    	  \begin{enumerate}
   		    			 \item[-] OptionSettingtype
   		    	\end{enumerate}
   	\item Contact List Layouts
	\item Contact Display Options
	\item Form for Email Options
	\item Integration Options
	\item Menuentry Options
 \end{itemize}
\\
\hline
\end{tabular}
\captionof{table}{Contact List All Category}
\end{minipage}

\begin{minipage}{0.7\textwidth}
\begin{tabular}{|p{0.6\textwidth} | p{0.6\textwidth}|}
\hline
\textbf{CSL} & \textbf{Joomla} \\ 
\hline
 1- List contact in a Category
   \begin{itemize}
     \item Basic  Setting 
    		\begin{enumerate}
    			\item[-] Select a Category
    		\end{enumerate}
    	\item Category Options
    	\item[+] Contacts in Category 
    		   \begin{enumerate}
    		    	 \item[-] OptionSettingtype
    		    \end{enumerate}
    \item Contact List Layouts
 	\item Contact Display Options
 	\item Form for Email Options
 	\item Integration Options
 	\item Menuentry Options
  \end{itemize}
 & 
1- List contact in a Category
  \begin{itemize}
    \item Basic  Setting 
   		\begin{enumerate}
   			\item[-] categorie\_id
   		\end{enumerate}
   	\item Category Options
   	\item[+] show\_cat\_items
   		    	  \begin{enumerate}
   		    			 \item[-] OptionSettingtype
   		    	\end{enumerate}
   	\item Contact List Layouts
	\item Contact Display Options
	\item Form for Email Options
	\item Integration Options
	\item Menuentry Options
 \end{itemize}
\\
\hline
\end{tabular}
\captionof{table}{List contact in a Category}
\end{minipage}

\begin{minipage}{0.7\textwidth}
\begin{tabular}{|p{0.6\textwidth} | p{0.6\textwidth}|}
\hline
\textbf{CSL} & \textbf{Joomla} \\ 
\hline
 1- List contact in a Category
   \begin{itemize}
     \item Basic  Setting 
    		\begin{enumerate}
    			\item[-] Select a Contact
    		\end{enumerate}
 	\item Contact Display Options
 	\item Form for Email Options
 	\item Menuentry Options
  \end{itemize}
 & 
1- Single Contact
  \begin{itemize}
    \item Basic  Setting 
   		\begin{enumerate}
   			\item[-] contact\_id
   		\end{enumerate}
	\item Contact Display Options
	\item Form for Email Options
	\item Menuentry Options
 \end{itemize}
\\
\hline
\end{tabular}
\captionof{table}{Single Contact}
\end{minipage}

\begin{minipage}{0.7\textwidth}
\begin{tabular}{|p{0.6\textwidth} | p{0.6\textwidth}|}
\hline
\textbf{CSL} & \textbf{Joomla} \\ 
\hline
 1- Featured Contact
 \begin{itemize}
    \item Contact List Layouts
 	\item Contact Display Options
 	\item Form for Email Options
 	\item Menuentry Options
  \end{itemize}
 & 
1- Featured Contact
 \begin{itemize}
   	\item Contact List Layouts
	\item Contact Display Options
	\item Form for Email Options
	\item Menuentry Options
 \end{itemize}
\\
\hline
\end{tabular}
\captionof{table}{Featured Contact}
\end{minipage}

\textbf{Web Links }
Wie schon bei anderen Pages bzw Menupunkt filtern wir erstmal  die wiederkehrenden Elementen aus dieser Menu-Categorie: Categories Options, Category Options, List Layout Options,und Menuentry Options.

\begin{minipage}{0.7\textwidth}
\begin{tabular}{|p{0.6\textwidth} | p{0.6\textwidth}|}
\hline
\textbf{CSL} & \textbf{Joomla} \\ 
\hline
\begin{itemize}
\item Display Select
 \begin{enumerate}
    \item[-] OptionSettingtype
  \end{enumerate} 
\item Table Headings
 \begin{enumerate}
    \item[-] OptionSettingtype
  \end{enumerate} 
\item Links description
 \begin{enumerate}
    \item[-] OptionSettingtype
  \end{enumerate} 
\item Hits
 \begin{enumerate}
    \item[-] OptionSettingtype
  \end{enumerate}  
\item Pagination
 \begin{enumerate}
    \item[-] OptionSettingtype
  \end{enumerate}  
\item Pagination Results
 \begin{enumerate}
    \item[-] OptionSettingtype
  \end{enumerate}  
\end{itemize} 
 & 
\begin{itemize}
\item show\_pagination\_limit
\begin{enumerate}
    \item[-] OptionSettingvalue
  \end{enumerate} 
\item show\_headings
\begin{enumerate}
    \item[-] OptionSettingvalue
  \end{enumerate} 
\item show\_link\_description
\begin{enumerate}
    \item[-] OptionSettingvalue
  \end{enumerate} 
\item show\_link\_hits
\begin{enumerate}
    \item[-] OptionSettingvalue
  \end{enumerate} 
\item show\_pagination
\begin{enumerate}
    \item[-] OptionSettingvalue
  \end{enumerate} 
\item show\_pagination\_results
\begin{enumerate}
    \item[-] OptionSettingvalue
  \end{enumerate} 
\end{itemize}
\\
\hline
\end{tabular}
\captionof{table}{Web Link List Layout}
\end{minipage}

\begin{minipage}{0.7\textwidth}
\begin{tabular}{|p{0.6\textwidth} | p{0.6\textwidth}|}
\hline
\textbf{CSL} & \textbf{Joomla} \\ 
\hline
 1- List All Web Link Categories
   \begin{itemize}
     \item Basic  Setting 
    		\begin{enumerate}
    			\item[-] Select a Category Level
    		\end{enumerate}
    	\item Categories Options
 	   	\begin{enumerate}
 	    \item[+] Top Level Category Description
 	    \item[+] Show /Hide the number of Web Links 
 	    	  \begin{enumerate}
 	    			 \item[-] OptionSettingtype
 	    	\end{enumerate}
 	   \end{enumerate}
    	\item Category Options
    	\item[+] Weblink in Category 
    		    	  \begin{enumerate}
    		    			 \item[-] OptionSettingtype
    		    	\end{enumerate}
    \item Weblink List Layouts
 	\item Menuentry Options
  \end{itemize}
 & 
1- List All Web Link Categories
  \begin{itemize}
    \item Basic  Setting 
   		\begin{enumerate}
   			\item[-] categorie\_id
   		\end{enumerate}
   	\item Categories Options
	   	\begin{enumerate}
	    \item[+]categories\_description
	    \item[+]show\_cat\_num\_ links\_ cat 
	    	  \begin{enumerate}
	    			 \item[-] OptionSettingtype
	    	\end{enumerate}
	   \end{enumerate}
   	\item Category Options
   	\item[+] show\_cat\_num\_links
   		    	  \begin{enumerate}
   		    			 \item[-] OptionSettingtype
   		    	\end{enumerate}
   	\item Weblinks List Layouts
	\item Menuentry Options
 \end{itemize}
\\
\hline
\end{tabular}
\captionof{table}{Weblinks List All Categories}
\end{minipage}

\begin{minipage}{0.7\textwidth}
\begin{tabular}{|p{0.6\textwidth} | p{0.6\textwidth}|}
\hline
\textbf{CSL} & \textbf{Joomla} \\ 
\hline
 1- List Web Link Category
   \begin{itemize}
     \item Basic  Setting 
    		\begin{enumerate}
    			\item[-] Select a Category 
    		\end{enumerate}
    	\item Category Options
    	\item[+] Weblink in Category 
    		    	  \begin{enumerate}
    		    			 \item[-] OptionSettingtype
    		    	\end{enumerate}
    \item Weblink List Layouts
    	\item Integration Options
 	\item Menuentry Options
  \end{itemize}
 & 
1- List Web Link Category
  \begin{itemize}
    \item Basic  Setting 
   		\begin{enumerate}
   			\item[-] categorie\_id
   		\end{enumerate}
   	\item Category Options
   	\item[+] show\_cat\_num\_links
   		    	  \begin{enumerate}
   		    			 \item[-] OptionSettingvalue
   		    	\end{enumerate}
   	\item Weblinks List Layouts
   	\item Integration Options
	\item Menuentry Options
 \end{itemize}
\\
\hline
\end{tabular}
\captionof{table}{List Web Link Category}
\end{minipage}

\begin{minipage}{0.7\textwidth}
\begin{tabular}{|p{0.6\textwidth} | p{0.6\textwidth}|}
\hline
\textbf{CSL} & \textbf{Joomla} \\ 
\hline
 1- Submit a Web Link
 \begin{itemize}
 	\item Menuentry Options
  \end{itemize}
 & 
1- Submit a Web Link
 \begin{itemize}
	\item Menuentry Options
 \end{itemize}
\\
\hline
\end{tabular}
\captionof{table}{Submit a Web Link}
\end{minipage}

\textbf{News Feed}

News Feed \{to do\} die wiederkehrende Elemente sind:

\begin{minipage}{0.63\textwidth}
\begin{tabular}{|p{0.6\textwidth} | p{0.6\textwidth}|}
\hline
\textbf{CSL} & \textbf{Joomla} \\ 
\hline
\begin{itemize}
\item Display Select
\begin{enumerate}
    \item[-] OptionSettingtype
 \end{enumerate} 
\item Table Headings
\begin{enumerate}
    \item[-] OptionSettingtype
 \end{enumerate} 
\item  Articles
\begin{enumerate}
    \item[-] OptionSettingtype
 \end{enumerate} 
\item Feed Links
\begin{enumerate}
    \item[-] OptionSettingtype
 \end{enumerate} 
\item Pagination
\begin{enumerate}
    \item[-] OptionSettingtype
 \end{enumerate} 
\item Pagination Results
\begin{enumerate}
    \item[-] OptionSettingtype
 \end{enumerate} 
\end{itemize}
 & 
\begin{itemize}
\item show\_pagination\_limit
 \begin{enumerate}
    \item[-] OptionSettingvalue
 \end{enumerate} 
\item show\_headings
 \begin{enumerate}
    \item[-] OptionSettingvalue
 \end{enumerate} 
\item show\_articles
 \begin{enumerate}
    \item[-] OptionSettingvalue
 \end{enumerate}  
\item show\_link
 \begin{enumerate}
    \item[-] OptionSettingvalue
 \end{enumerate} 
\item show\_pagination
 \begin{enumerate}
    \item[-] OptionSettingvalue
 \end{enumerate}  
\item show\_pagination\_results 
\begin{enumerate}
    \item[-] OptionSettingvalue
 \end{enumerate}  
\end{itemize}
\\
\hline
\end{tabular}
\captionof{table}{News Feeds List Layouts}
\end{minipage}

\begin{minipage}{0.63\textwidth}
\begin{tabular}{|p{0.55\textwidth} | p{0.57\textwidth}|}
\hline
\textbf{CSL} & \textbf{Joomla} \\ 
\hline
\begin{itemize}
\item Feed Image
\begin{enumerate}
    \item[-] OptionSettingtype
 \end{enumerate}  
\item Feed Description
\begin{enumerate}
    \item[-] OptionSettingtype
 \end{enumerate}  
\item Feed Content
\begin{enumerate}
    \item[-] OptionSettingtype
 \end{enumerate}  
\item Characters count 
 \end{itemize}
 &
 \begin{itemize}
\item show\_feed\_image
\begin{enumerate}
    \item[-] OptionSettingvalue
 \end{enumerate} 
\item show\_feed\_description
\begin{enumerate}
    \item[-] OptionSettingvalue
 \end{enumerate}  
\item show\_item\_description
\begin{enumerate}
    \item[-] OptionSettingvalue
 \end{enumerate}  
\item feed\_character\_count 
 \end{itemize} 
\\
\hline
\end{tabular}
\captionof{table}{Feed Display Options}
\end{minipage}

\begin{minipage}{0.7\textwidth}
\begin{tabular}{|p{0.6\textwidth} | p{0.6\textwidth}|}
\hline
\textbf{CSL} & \textbf{Joomla} \\ 
\hline
 1-List All News Feeds Categories
   \begin{itemize}
     \item Basic  Setting 
    		\begin{enumerate}
    			\item[-] Select a Category Level
    		\end{enumerate}
    	\item Categories Options
 	   	\begin{enumerate}
 	    \item[+] Top Level Category Description
 	    \item[+] Show /Hide the number of News Feeds 
 	    	  \begin{enumerate}
 	    			 \item[-] OptionSettingtype
 	    	\end{enumerate}
 	   \end{enumerate}
    	\item Category Options
    	\item[+] Feeds in Category 
    		    	  \begin{enumerate}
    		    			 \item[-] OptionSettingtype
    		    	\end{enumerate}
    \item News Feeds List Layouts
    \item Feed Display Options
 	\item Menuentry Options
  \end{itemize}
 & 
1- List All News Feeds Categories
  \begin{itemize}
    \item Basic  Setting 
   		\begin{enumerate}
   			\item[-] categorie\_id
   		\end{enumerate}
   	\item Categories Options
	   	\begin{enumerate}
	    \item[+]categories\_description
	    \item[+]show\_cat\_items\_cat 
	    	  \begin{enumerate}
	    			 \item[-] OptionSettingtype
	    	\end{enumerate}
	   \end{enumerate}
   	\item Category Options
   	\item[+] show\_cat\_items
   		    	  \begin{enumerate}
   		    			 \item[-] OptionSettingtype
   		    	\end{enumerate}
   	\item Weblinks List Layouts
   	\item News Feeds List Layouts
   	\item Feed Display Options
	\item Menuentry Options
 \end{itemize}
\\
\hline
\end{tabular}
\captionof{table}{List All News Feeds Categories}
\end{minipage}

\begin{minipage}{0.7\textwidth}
\begin{tabular}{|p{0.6\textwidth} | p{0.6\textwidth}|}
\hline
\textbf{CSL} & \textbf{Joomla} \\ 
\hline
 1-List news Feeds in a Category
   \begin{itemize}
     \item Basic  Setting 
    		\begin{enumerate}
    			\item[-] Select a Category 
    		\end{enumerate}
    	\item Category Options
    	\item[+] Feeds in Category 
    		    	  \begin{enumerate}
    		    			 \item[-] OptionSettingtype
    		    	\end{enumerate}
    \item News Feeds List Layouts
    \item Feed Display Options
    \item[+] Feed display order
           	\begin{enumerate}
           		\item[-] Use Global
           		\item[-] Most recent first
           		\item[-] Oldest first
           	\end{enumerate}	
 	\item Menuentry Options
  \end{itemize}
 & 
1- List news Feeds in a Category
  \begin{itemize}
    \item Basic  Setting 
   		\begin{enumerate}
   			\item[-] categorie\_id
   		\end{enumerate}
   	\item Category Options
   	\item[+] show\_cat\_items
		\begin{enumerate}
 			 \item[-] OptionSettingvalue
	    	\end{enumerate}
   	\item News Feeds List Layouts
   	\item Feed Display Options
   	\item[+] feed\_display\_order
   	   	\begin{enumerate}
   	   		\item[-] leer
   	   		\item[-] desc
   	   		\item[-] asc
   	   	\end{enumerate}	
	\item Menuentry Options
 \end{itemize}
\\
\hline
\end{tabular}
\captionof{table}{List news Feeds in a Category}
\end{minipage}

\begin{minipage}{0.7\textwidth}
\begin{tabular}{|p{0.6\textwidth} | p{0.6\textwidth}|}
\hline
\textbf{CSL} & \textbf{Joomla} \\ 
\hline
 1-List news Feeds in a Category
   \begin{itemize}
     \item Basic  Setting 
    		\begin{enumerate}
    			\item[-] Select a Feed
    		\end{enumerate}  
    \item Feed Display Options
    \item[+] Feed display order
           	\begin{enumerate}
           		\item[-] Use Global
           		\item[-] Most recent first
           		\item[-] Oldest first
           	\end{enumerate}	
 	\item Menuentry Options
  \end{itemize}
 & 
1- List news Feeds in a Category
  \begin{itemize}
    \item Basic  Setting 
   		\begin{enumerate}
   			\item[-] feed\_id
   		\end{enumerate}
   	\item Feed Display Options
   		\item[+] feed\_display\_order
   	   	\begin{enumerate}
   	   		\item[-] leer
   	   		\item[-] desc
   	   		\item[-] asc
   	   	\end{enumerate}	
	\item Menuentry Options
 \end{itemize}
\\
\hline
\end{tabular}
\captionof{table}{Single News Feed}
\end{minipage}

\textbf{Smart Search}

diese Page ist \{to do\}

\begin{minipage}{0.7\textwidth}
\begin{tabular}{|p{0.6\textwidth} | p{0.6\textwidth}|}
\hline
\textbf{CSL} & \textbf{Joomla} \\ 
\hline
\begin{itemize}
\item Search Query
\item Search Filter
 \begin{enumerate}
 \item[-]  Search Filter
 \end{enumerate}
\item Date Filters
\begin{enumerate}
 	\item[-] OptionSettingtype
 \end{enumerate} 
\item Advanced Search
\begin{enumerate}
 	\item[-] OptionSettingtype
 \end{enumerate} 
\item Expand Advanced Search
\begin{enumerate}
 	\item[-] OptionSettingtype
 \end{enumerate} 
\item Result Description
\begin{enumerate}
 	\item[-] OptionSettingtype
 \end{enumerate} 
\item Description Length 
\item Result URL
\begin{enumerate}
 	\item[-] OptionSettingtype
 \end{enumerate}
\item Display Select
\begin{enumerate}
 	\item[-] OptionSettingvalue
 \end{enumerate} 
\end{itemize}

 & 
\begin{itemize}
\item q
\item f
 \begin{enumerate}
 \item[-] filter\_id
 \end{enumerate}
\item show\_date\_filters
\begin{enumerate}
 	\item[-] OptionSettingvalue
 \end{enumerate} 
\item show\_advanced
\begin{enumerate}
 	\item[-] OptionSettingvalue
 \end{enumerate} 
\item expand\_advanced
\begin{enumerate}
 	\item[-] OptionSettingvalue
 \end{enumerate} 
\item show\_description
\begin{enumerate}
 	\item[-] OptionSettingvalue
 \end{enumerate} 
\item description\_length 
\item show\_url
\begin{enumerate}
 	\item[-] OptionSettingvalue
 \end{enumerate}
\item show\_pagination\_limit
\begin{enumerate}
 	\item[-] OptionSettingvalue
 \end{enumerate} 
\end{itemize}
\\
\hline
\end{tabular}
\captionof{table}{MaxlevelSettingtype}
\end{minipage}

\begin{minipage}{0.7\textwidth}
\begin{tabular}{|p{0.6\textwidth} | p{0.6\textwidth}|}
\hline
\textbf{CSL} & \textbf{Joomla} \\ 
\hline
\begin{itemize}
\item Pagination
\begin{enumerate}
 	\item[-] OptionSettingvalue
 \end{enumerate}
\item Pagination Results
\begin{enumerate}
 	\item[-] OptionSettingvalue
 \end{enumerate} 
\item Allow Empty Search
\begin{enumerate}
 	\item[-] No
 	\item[-] yes 
 \end{enumerate} 
\item Sort Field
\begin{enumerate}
 	\item[-] Use Global
 	\item[-] relevance
 	\item[-] date
 	\item[-] price
 \end{enumerate} 
\item Sort Direction
\begin{enumerate}
 	\item[-] Use Global
 	\item[-] descending
 	\item[-] ascending
 \end{enumerate} 
\item show\_feed
\begin{enumerate}
 	\item[-] No
 	\item[-] Yes
 \end{enumerate} 
\item Show Feed
\begin{enumerate}
 	\item[-] No
 	\item[-] Yes
 \end{enumerate} 
\item Show Feed Text 
\item Menuentry Options
\end{itemize}
 & 
\begin{itemize}
\item show\_pagination
\begin{enumerate}
 	\item[-] OptionSettingvalue
 \end{enumerate}
\item show\_pagination\_results
\begin{enumerate}
 	\item[-] OptionSettingvalue
 \end{enumerate} 
\item allow\_empty\_query
\begin{enumerate}
 	\item[-] 0
 	\item[-] 1
 \end{enumerate} 
\item sort\_order
\begin{enumerate}
 	\item[-] leer
 	\item[-] relevance
 	\item[-] date
 	\item[-] price
 \end{enumerate} 
\item sort\_direction
\begin{enumerate}
 	\item[-] leer
 	\item[-] desc
 	\item[-] asc
 \end{enumerate} 
\item show\_feed
\begin{enumerate}
 	\item[-] 0
 	\item[-] 1
 \end{enumerate} 
\item show\_feed\_text
\begin{enumerate}
 	\item[-] 0
 	\item[-] 1
 \end{enumerate} 
\item show\_feed\_link 
\item Menuentry Options
\end{itemize}
\\
\hline
\end{tabular}
\captionof{table}{Search}
\end{minipage}

\textbf{Search}

\begin{minipage}{0.7\textwidth}
\begin{tabular}{|p{0.6\textwidth} | p{0.6\textwidth}|}
\hline
\textbf{CSL} & \textbf{Joomla} \\ 
\hline
\begin{itemize}
 \item Search Term 
 \item Use Search Areas
 \begin{enumerate}
  	\item[-] OptionSettingtype
  \end{enumerate} 
 \item Created Date
 \begin{enumerate}
  	\item[-] OptionSettingtype
  \end{enumerate}
 \item Search for
 \begin{enumerate}
  	\item[-] All Words
  	\item[-] Any Words
  	\item[-] Exact Phrase
  \end{enumerate} 
 \item Results ordering
 \begin{enumerate}
  	\item[-] Newest First
  	\item[-] Oldest First
  	\item[-] Popularity
  	\item[-] Alphabetical
  	\item[-] Category
  \end{enumerate} 
  \item Menuentry Options
 \end{itemize}

 & 
 \begin{itemize}
 \item searchword
 \item search\_areas
 \begin{enumerate}
  	\item[-] OptionSettingvalue
  \end{enumerate} 
 \item show\_date
 \begin{enumerate}
  	\item[-] OptionSettingvalue
  \end{enumerate}
 \item searchphrase
 \begin{enumerate}
  	\item[-] 0
  	\item[-] 1
  	\item[-] 2
  \end{enumerate} 
 \item ordering
 \begin{enumerate}
  	\item[-] newest
  	\item[-] oldest
  	\item[-] popular
  	\item[-] alpha
  	\item[-] category
  \end{enumerate} 
  \item Menuentry Options
 \end{itemize}
\\
\hline
\end{tabular}
\captionof{table}{Search}
\end{minipage}

\textbf{Iframe Wrapper}

\begin{minipage}{0.7\textwidth}
\begin{tabular}{|p{0.6\textwidth} | p{0.6\textwidth}|}
\hline
\textbf{CSL} & \textbf{Joomla} \\ 
\hline
\begin{itemize}
\item URL 
\item Scroll Bars
\begin{enumerate}
	\item[-] no
	\item[-] yes
	\item[-] auto
\end{enumerate} 
\item Width 
\item Height 
\item Auto height
\begin{enumerate}
	\item[-] YesNoSettingtype
\end{enumerate} 
\item Auto add
\begin{enumerate}
	\item[-] YesNoSettingtype
\end{enumerate}
\item Frame border
\begin{enumerate}
	\item[-] YesNoSettingtype
\end{enumerate}
\item Menuentry Options
\end{itemize}
 & 
\begin{itemize}
\item url 
\item scrolling
\begin{enumerate}
	\item[-] no
	\item[-] yes
	\item[-] auto
\end{enumerate} 
\item width 
\item height 
\item height\_auto
\begin{enumerate}
	\item[-] YesNoSettingvalue
\end{enumerate} 
\item add\_scheme
\begin{enumerate}
	\item[-] YesNoSettingvalue
\end{enumerate}
\item frameborder
\begin{enumerate}
	\item[-] YesNoSettingvalue
\end{enumerate}
\item Menuentry Options
\end{itemize}
\\
\hline
\end{tabular}
\captionof{table}{Iframe Wrapper}
\end{minipage}

\textbf{External URL}

\begin{minipage}{0.7\textwidth}
\begin{tabular}{|p{0.6\textwidth} | p{0.6\textwidth}|}
\hline
\textbf{CSL} & \textbf{Joomla} \\ 
\hline
1-External URL
 \begin{itemize}
 \item Link Type Options
 \end{itemize}
 & 
1-External URL
\begin{itemize}
 \item Link Type Options
 \end{itemize}
\\
\hline
\end{tabular}
\captionof{table}{External URL}
\end{minipage}

\textbf{Menu Item Alias}

\begin{minipage}{0.7\textwidth}
\begin{tabular}{|p{0.6\textwidth} | p{0.6\textwidth}|}
\hline
\textbf{CSL} & \textbf{Joomla} \\ 
\hline

 \begin{itemize}
 \item Select a Menu Item
 \item Link Type Options
 \end{itemize}
 & 

\begin{itemize}
\item itemid
 \item Link Type Options
 \end{itemize}
\\
\hline
\end{tabular}
\captionof{table}{Menu Item Alias}
\end{minipage}

\textbf{System Link-Text Separator}

\begin{minipage}{0.7\textwidth}
\begin{tabular}{|p{0.6\textwidth} | p{0.6\textwidth}|}
\hline
\textbf{CSL} & \textbf{Joomla} \\ 
\hline

 \begin{itemize}
 \item Link Type Options
 \end{itemize}
 & 

\begin{itemize}
 \item Link Type Options
 \end{itemize}
\\
\hline
\end{tabular}
\captionof{table}{System Link-Text Separator}
\end{minipage}

\chapter{Anhang 2}



%% Dokument ENDE %%%%%%%%%%%%%%%%%%%%%%%%%%%%%%%%%%%%%%%%%%%%%%%%%%%%%%%%%%
\end{document}

